\documentclass[8pt,landscape]{extarticle}
\usepackage[letterpaper,landscape,margin=0.18in,top=0.18in,bottom=0.18in]{geometry}
\usepackage{multicol}
\usepackage{amsmath,amssymb}
\usepackage{array}
\usepackage{booktabs}
\usepackage{xcolor}
\usepackage{listings}
\usepackage{enumitem}
\usepackage{titlesec}
\usepackage[T1]{fontenc}
\usepackage{lmodern}
\usepackage{microtype}

\AtBeginDocument{\small}

% ---------- Spacing / layout ----------
\setlength{\columnsep}{0.3em}
\setlength{\columnseprule}{0.15pt}
\setlength{\parindent}{0pt}
\setlength{\parskip}{1pt}
\setlist{nosep,leftmargin=0.9em,itemsep=0pt,topsep=0.5pt,parsep=0pt,partopsep=0pt}
\setlength{\tabcolsep}{2pt}
\renewcommand{\arraystretch}{1.0}
\pagestyle{empty}
\raggedcolumns

% ---------- Section styling (tight) ----------
\titleformat{\section}{\normalfont\bfseries\footnotesize\color{blue!60!black}}
  {\thesection.}{0.25em}{}
\titlespacing*{\section}{0pt}{3pt}{0.5pt}
\titleformat{\subsection}{\normalfont\bfseries\scriptsize}{}{0em}{}
\titlespacing*{\subsection}{0pt}{2pt}{0.2pt}

% ---------- ARM listings ----------
\lstdefinelanguage{arm}{
  morekeywords={ldr,ldrb,ldrh,ldrsb,ldrsh,ldrd,str,strb,strh,strd,mov,movw,movt,%
    movs,add,adds,adc,sub,subs,sbc,cmp,cmn,tst,teq,bl,bx,b,beq,bne,bgt,bge,blt,%
    ble,bhi,bhs,blo,bls,bcs,bcc,push,pop,cbz,cbnz,it,itt,ite,itte,itet,itee,%
    ittt,ittee,iteee,ittte,itttt,tbb,tbh,adr,lsl,lsls,lsr,lsrs,asr,asrs,ror,%
    rors,rrx,rrxs,and,ands,orr,orrs,eor,eors,bic,bics,orn,mvn,mvns,bfc,bfi,%
    mul,muls,mla,mls,umull,umlal,smull,smlal,udiv,sdiv,nop,wfi,wfe,svc,rsb,%
    lslge,lslgt,lsllt,lslle,asrge,asrgt,asrlt,asrle,ldrle,ldrgt,ldrlt,ldrge,%
    ldrshgt,ldrshlt,ldrhge,ldrhlt,lsrlt,lsrge},
  morecomment=[l]{@},
  sensitive=false
}
\lstset{
  language=arm,
  basicstyle=\ttfamily\scriptsize,
  keywordstyle=\color{blue!70!black}\bfseries,
  commentstyle=\color{green!40!black}\itshape,
  breaklines=false,
  frame=none,
  aboveskip=0.5pt,
  belowskip=0.5pt,
  columns=fullflexible,
  keepspaces=true,
  xleftmargin=0.25em,
  showstringspaces=false
}
\newcommand{\asm}[1]{\texttt{\scriptsize #1}}
\newcommand{\cee}[1]{\texttt{\scriptsize #1}}

% Example / Problem tags
\newenvironment{ex}
  {\par\smallskip\noindent{\scriptsize\color{red!55!black}\textbf{Ex.}}\,\scriptsize\ignorespaces}
  {\par}
\newenvironment{qp}
  {\par\smallskip\noindent{\scriptsize\color{purple!70!black}\textbf{Qz.}}\,\scriptsize\ignorespaces}
  {\par}
\newenvironment{hwp}
  {\par\smallskip\noindent{\scriptsize\color{orange!70!black}\textbf{HW}}\,\scriptsize\ignorespaces}
  {\par}

\begin{document}
\begin{center}
  \textbf{\footnotesize CEC 320 --- FINAL EXAM Cheatsheet} \;
  \scriptsize HW 5--12 \textbullet{} Lctr 14--27 \textbullet{} Quiz 3, 4, 6a
\end{center}
\vspace{-8pt}

\begin{multicols*}{4}
\scriptsize

% ==========================================================
% ==========  PART 0: LECTURE 14--27 RECALL BULLETS  =======
% ==========================================================

\section{Lctr 14--27 Recall Bullets}

\subsection{L14 Mixed C/ASM}
\begin{itemize}
\item EABI: R0--R3 args+ret, caller-saved. R4--R11 callee-saved.
\item 64-bit arg/ret: low in R0, high in R1.
\item Stack = Full-Descending: SP at last used; PUSH decrements then writes.
\item Multi-push: low regs at low addrs regardless of brace order.
\item Return: \asm{bx lr} (no LR push) or \asm{pop \{pc\}} (LR pushed).
\end{itemize}

\subsection{L15 Arithmetic}
\begin{itemize}
\item \asm{ADD/SUB/RSB} (no S) $\not\to$ flags. \asm{ADDS/SUBS} $\to$ NZCV.
\item \asm{ADC Rn,Op2}: Rn+Op2+C. \asm{SBC Rn,Op2}: Rn$-$Op2$-$(1$-$C).
\item 64-bit add: \asm{adds/adc}. 64-bit sub: \asm{subs/sbc}.
\item \asm{MUL/MLA/MLS}: 32-bit. \asm{SMULL/UMULL/SMLAL/UMLAL}: 64-bit.
\item Modulo: \asm{udiv/mul/sub} or \asm{udiv} then \asm{mls Rd,Rn,Rq,Ra}.
\end{itemize}

\subsection{L16 Shifts / Op2}
\begin{itemize}
\item LSL s\&u. LSR u-only (0 fill). ASR s-only (sign repeat). S-suf $\to$ C=last bit out.
\item Op2 forms: \#imm / Rm / Rm,\{LSL/LSR/ASR/ROR\} \#n.
\item Shifted-Op2: dest can NOT be omitted.
\item Q31 mul (3-instr): \asm{smull r1,r0,r0,r1 / lsl r0,\#1 / add r0,r0,r1,lsr \#31}.
\item Fast-scale by $(2^n\pm1)/2^n$: \asm{add r0,r0,r0,asr \#n} or \asm{sub ... asr \#n}.
\end{itemize}

\subsection{L17 Bitwise}
\begin{itemize}
\item AND/ORR/EOR/MVN. BIC = \cee{Rn \& \textasciitilde Op2}. ORN = \cee{Rn | \textasciitilde Op2}. No C op for either.
\item BFC \asm{Rd,\#s,\#w} clears w bits at s. BFI inserts w LSBs.
\item APSR: \asm{MRS Rd,apsr} read; \asm{MSR apsr\_nzcvq,Rn} write.
\item Bitfield assign: BIC mask, ORR value-shifted-in.
\end{itemize}

\subsection{L18 Flags}
\begin{itemize}
\item N=MSB, Z=(result==0), C=unsigned carry/!borrow, V=signed ovf.
\item HW updates BOTH C and V every arithmetic; user picks.
\item SUB: C=1 means NO borrow (\cee{u1{\mbox{>}}=u2}), C=0 means borrow.
\item N-bit sim trick: shift operands left $32{-}N$, do ADDS, shift back ASR, read APSR.
\end{itemize}

\subsection{L19 CCS}
\begin{itemize}
\item \asm{CMP}=SUB no-writeback. \asm{CMN}=ADD no-writeback.
\item \asm{cmp r0,\#-5} $\Rightarrow$ assembler emits \asm{CMN r0,\#5}.
\item Signed CCS: GT/GE/LT/LE. Unsigned: HI/HS/LO/LS. HS=CS, LO=CC.
\item Match type to CCS: signed ops $\to$ signed CCS + ASR; unsigned $\to$ unsigned + LSR.
\end{itemize}

\subsection{L20 LDR/STR (imm)}
\begin{itemize}
\item \asm{[Ra,\#os]} basic; \asm{[Ra,\#os]!} update BEFORE; \asm{[Ra],\#os} update AFTER.
\item LOAD sufs: ---/B/SB/H/SH. STORE: ---/B/H (no sign).
\item C$\leftrightarrow$ASM: \cee{*p++}$=$post; \cee{*++p}$=$pre; \cee{*(p+k)}$=$basic.
\item STRB writes only low byte of Rv; upper bits ignored.
\item Offset ranges: word $\pm$4095; byte/half $\pm$255 on some forms.
\end{itemize}

\subsection{L21 LDR/STR (reg) + MOV}
\begin{itemize}
\item \asm{[Ra,Ri,LSL \#n]} with $n=\log_2$(elem size). No writeback form.
\item MOV small imm. MOVW low16. MOVT high16.
\item Pseudo \asm{ldr Rd,=expr} picks MOV/MOVW/literal-pool by size.
\item Two array iters: post-inc ptr vs fixed-base + reg-offset.
\end{itemize}

\subsection{L22 Inc/Dec}
\begin{itemize}
\item C precedence: postfix \cee{++} L1 (LTR); prefix, \cee{*}, \cee{\&}, cast L2 (RTL).
\item \cee{*p++}=deref old, advance. \cee{*++p}=advance, deref new.
\item \cee{++*p} = val++ in mem (ldr/add/str). \cee{++*p++} = post-inc p, store-back via \asm{[r0,\#-1]}.
\item Ptr arith scaled: p8+1, p16+2, p32+4. Typed diff $\to$ elements.
\end{itemize}

\subsection{L23 ASM Calls ASM}
\begin{itemize}
\item Leaf: no BL $\to$ \asm{bx lr}, no LR push.
\item Stem: has BL $\to$ \asm{push \{lr\}} + \asm{pop \{pc\}}.
\item LDRD \asm{Rv\_L,Rv\_H,[Ra,\#os]}: 8 bytes, L=low addr, little-endian.
\item No reg-offset form for LDRD/STRD.
\end{itemize}

\subsection{L24 Combo Branch / CEX}
\begin{itemize}
\item \asm{CBZ/CBNZ Rn,lbl}: fwd only ($\le$126B), Rn=R0--R7, no flags.
\item Perfect after LDRB (byte value already in reg, null-term test).
\item IT/ITT/ITE/ITTT/ITTE/ITEE/ITTEE/ITTTT --- up to 4 slots.
\item First=T=cond; later T=cond, E=inverse.
\item CBZ/CBNZ FORBIDDEN in IT. Cond B allowed ONLY as LAST slot.
\end{itemize}

\subsection{L25 If/If-else}
\begin{itemize}
\item PL: same-CCS branch to \_then; else FIRST.
\item NL: inverse-CCS branch to \_else; then FIRST. Reads like C.
\item if-elseif: chain CMP+NL-branch, \asm{b end} each arm.
\item Compound OR: double CMP $\to$ shared \_then; fall-through=else.
\item Compound AND (De Morgan): both CMPs $\to$ \_else; fall-through=then.
\end{itemize}

\subsection{L26 Loops}
\begin{itemize}
\item Labels \_lp (top), \_end (exit). while/for top-tested (0+). do-while bottom-tested ($\ge$1).
\item while: NL cmp at top to \_end, body, \asm{b lp}.
\item while-late: jump over body on entry, test at bottom.
\item while-CEX (rare): \asm{cmp/it cond/bcond body\_in\_it/b end}.
\item Priming-read duplicates LDRB; while(1)+break collapses.
\end{itemize}

\subsection{L27 break/continue/switch}
\begin{itemize}
\item break = fwd \asm{b end}. continue (while) = back to cond.
\item continue in \texttt{for} = branch to UPDATE expr (i++ still runs).
\item Switch naive: CMP/BEQ ladder; stacked labels = fall-through.
\item \asm{ADR Rd,lbl} PC-rel addr; \asm{TBB [Rn,Rm]}: PC=PC+4+2$\cdot$table[Rm].
\item TBB table bytes: \texttt{(case\_X-first\_case)/2}.
\end{itemize}

% ==========================================================
% ==========  PART I: ARITHMETIC, FLAGS, BITWISE  ==========
% ==========================================================

\section{Q-Format Fixed-Point (HW5, L14)}
\textbf{UQm.n} unsigned, $N=m+n$. \textbf{Qm.n} signed, $N=1{+}m{+}n$.

\textbf{Encode:} $I = \mathrm{round}(f \cdot 2^n)$; if $<0$ add $2^N$, then hex.\\
\textbf{Decode:} if MSB=1 subtract $2^N$ (two's comp); $f = I/2^n$.

\begin{tabular}{@{}ll@{}}
\toprule
Format & Scale \\
\midrule
Q15 / Q1.15   & $2^{15}=32768$ \\
Q15.16        & $2^{16}=65536$ \\
Q16.15        & $2^{15}=32768$ (32-bit signed) \\
Q3.4 (8-bit)  & $2^{4}=16$ \\
UQ3.5 (8-bit) & $2^{5}=32$ \\
\bottomrule
\end{tabular}

\begin{hwp}\textbf{HW5 P4:} encode $6.73$ in Q3.4 $\to$ $6.73\cdot16=107.68\to108=$\textbf{0x6C}; encode $-4.5$ $\to$ $-72+256=184=$\textbf{0xB8}.\end{hwp}

\begin{hwp}\textbf{HW5 P5 Q15 MAC:} $f_1f_2+f_3f_4$: encode each ($0.5{\to}$0x4000, $0.25{\to}$0x2000, $0.125{\to}$0x1000, $0.0625{\to}$0x0800). Product: $(I_A{\cdot}I_B)\mathbin{>}\mathbin{>}15$. Sum results (plain int add).\end{hwp}

\textbf{Multiply rule:} Qm.n $\times$ Qp.q $=$ Qm+p+1.n+q (in 2N-bit). Shift right by $n$ (or $q$) to renormalize to Qm.n (or Qp.q).

\textbf{HW7 P1 Q15.16 extraction:} after \asm{smull r1,r0,r0,r1} (R0=HW, R1=LW), result bits b47..b16:
\begin{lstlisting}
smull r1,r0,r0,r1
lsl   r0,#16          @ HW << (47-31)
add   r0,r0,r1,lsr #16
\end{lstlisting}
Rule: $\mathrm{lsl}$ HW by (target\_MSB$-31$); $\mathrm{lsr}$ LW by $n$. Q31: lsl1/lsr31. Q15.16: lsl16/lsr16. Q16.15: lsl17/lsr15.

% ==========================================================
\section{Arithmetic Instrs (L14--15)}
\textbf{S-suffix rule:} only S-variants update N,Z,C,V. \asm{ADD} does NOT; \asm{ADDS} does. \asm{MULS} sets only N,Z (no C/V). \asm{UDIV/SDIV} NEVER set flags (no S form).

\subsection{ADDS/ADC --- 64-bit add}
\begin{lstlisting}
@ R0=lo(A), R1=hi(A), R2=lo(B), R3=hi(B)
mp_add_64bit_s:
  adds r0,r2    @ lo, sets C
  adc  r1,r3    @ hi + C
  bx   lr
\end{lstlisting}
\begin{hwp}\textbf{HW6 P5:} 0x88{\textunderscore}000000+0x89{\textunderscore}000000: lo=0x11{\textunderscore}000000 \textbf{C=1}; hi=3+2+1=6.\end{hwp}

\subsection{SUBS/SBC --- 64-bit sub}
\begin{lstlisting}
mp_sub_64bit_s:
  subs r0,r2    @ lo, sets C=NOT borrow
  sbc  r1,r3    @ hi - r3 - (1-C)
  bx   lr
\end{lstlisting}
\textbf{ARM rule:} C=1 $\Rightarrow$ NO borrow (i.e.\ $u_1\ge u_2$); C=0 $\Rightarrow$ borrow. Opposite of x86.

\subsection{Modulo via UDIV/MUL}
$A\%B = A - (A/B){\cdot}B$. Leaf, no push/pop:
\begin{lstlisting}
mp_umod_32bit_s:
  udiv r2,r0,r1
  mul  r2,r2,r1
  sub  r0,r0,r2
  bx   lr
\end{lstlisting}
MUL alone (no S) never sets flags.

\subsection{Signed/unsigned 64-bit products}
\begin{tabular}{@{}ll@{}}
\toprule
Insn & Operation \\
\midrule
\asm{UMULL Rlo,Rhi,Rn,Rm}   & unsigned $\{Rhi{:}Rlo\}=Rn{\cdot}Rm$ \\
\asm{SMULL Rlo,Rhi,Rn,Rm}   & signed $\{Rhi{:}Rlo\}=Rn{\cdot}Rm$ \\
\asm{UMLAL Rlo,Rhi,Rn,Rm}   & unsigned $\{Rhi{:}Rlo\}\mathbin{+}{=}Rn{\cdot}Rm$ \\
\asm{SMLAL Rlo,Rhi,Rn,Rm}   & signed $\{Rhi{:}Rlo\}\mathbin{+}{=}Rn{\cdot}Rm$ \\
\asm{MLA Rd,Rn,Rm,Ra}       & $Rd=Ra+Rn{\cdot}Rm$ (32-bit) \\
\asm{MLS Rd,Rn,Rm,Ra}       & $Rd=Ra-Rn{\cdot}Rm$ (32-bit) \\
\bottomrule
\end{tabular}

\begin{hwp}\textbf{HW8 P4} (sum-of-squares, 64b):
\begin{lstlisting}
  ldr r2,=0
  ldr r3,=0
lp:
  cmp   r0,r1
  bgt   end
  smlal r2,r3,r0,r0 @ {r3:r2}+=i*i
  add   r0,#1
  b     lp
end:
  mov r0,r2
  mov r1,r3
  bx  lr
\end{lstlisting}
Return: R0=lo, R1=hi per AAPCS.\end{hwp}

% ==========================================================
\section{EABI Arg Passing (L15, HW6)}
First 4 args $\to$ R0--R3, then stack. Sub-word args are promoted to 32-bit AT THE CALL SITE:
\begin{itemize}
\item \cee{uint8\_t 0xF} $\to$ R0 = 0x0000000F (zero-ext)
\item \cee{uint8\_t 0xFF} $\to$ R1 = 0x000000FF (zero-ext)
\item \cee{int8\_t 0x81} ($=-127$) $\to$ R1 = \textbf{0xFFFFFF81} (sign-ext)
\item \cee{int8\_t 0x7E} ($=+126$) $\to$ R0 = 0x0000007E
\end{itemize}

\textbf{Return:} 32-bit in R0. 64-bit in R0:R1 (R0=lo). Sub-word returns go in R0 (implicitly extended by caller).

\textbf{Q-format return:} encode as described above, put hex into R0. Q15.16 of $3.25$: $I = 3.25 {\cdot} 2^{16} =$ \textbf{0x00034000}.

% ==========================================================
\section{Stack \& PUSH/POP (L15)}
ARM uses \textbf{full-descending} stack: SP points at top-most used word; PUSH decrements SP then writes.

\textbf{PUSH ordering rule:} Regardless of list syntax, \asm{push \{list\}} writes lowest reg\# at LOWEST address (ascending). Listing order is irrelevant.

\begin{hwp}\textbf{HW6 P1:} SP=0x20001020; \asm{push \{r7,r5,r8,r6\}} $\to$ SP=0x20001010; storage:
\begin{tabular}{@{}ll@{}}
0x..1010 & R5 \\
0x..1014 & R6 \\
0x..1018 & R7 \\
0x..101C & R8 \\
\end{tabular}
\end{hwp}

\begin{qp}\textbf{Q3 P3:} SP=0x20001020; \asm{push \{r4,r8,r6\}} $\to$ SP=0x20001014; R4@1014, R6@1018, R8@101C. R6=(1{\mbox{<}}{\mbox{<}}6)+(1{\mbox{<}}{\mbox{<}}8)=\textbf{0x00000140}; byte@1019 = second byte of R6 = 0x01 (little-endian).\end{qp}

\asm{push \{r4-r7,lr\}}: SP$-$=20; sorted low$\to$hi. \asm{pop \{r4-r7,pc\}}: restore and return (folds return into pop).

% ==========================================================
\section{Barrel Shifter Op2 (L15--16)}
Op2 forms: (a) \#imm; (b) \asm{Rm}; (c) \asm{Rm, \{LSL|LSR|ASR|ROR\} \#n}. Shift count MUST be imm in Op2 (standalone LSL/LSR can use Rm for count).

\textbf{GNU quirk:} with shifted-reg Op2, dest register cannot be omitted --- \asm{add r0,r1,lsl \#3} is illegal.

\subsection{Shifted-Op2 scaling idioms}
\begin{tabular}{@{}ll@{}}
\toprule
ASM & C equiv \\
\midrule
\asm{lsl r0,r1,\#n}                & $2^n{\cdot}r_1$ (s or u) \\
\asm{add r0,r1,r1,lsl \#1}         & $3{\cdot}r_1$ \\
\asm{add r0,r1,r1,lsl \#2}         & $5{\cdot}r_1$ \\
\asm{rsb r0,r1,r1,lsl \#3}         & $7{\cdot}r_1$ \\
\asm{add r0,r1,r1,lsl \#4}         & $17{\cdot}r_1$ \\
\asm{rsb r0,r1,r1,lsl \#4}         & $15{\cdot}r_1$ \\
\asm{asr r0,r1,\#n}                & signed $r_1/2^n$ \\
\asm{lsr r0,r1,\#n}                & unsigned $r_1/2^n$ \\
\asm{add r0,r1,r1,asr \#2}         & $1.25{\cdot}r_1$ signed \\
\asm{sub r0,r1,r1,asr \#2}         & $0.75{\cdot}r_1$ signed \\
\asm{add r0,r1,r1,asr \#3}         & $1.125{\cdot}r_1$ signed \\
\asm{sub r0,r1,r1,asr \#3}         & $0.875{\cdot}r_1$ signed \\
\asm{rsb r0,r1,r1,asr \#2}         & $-0.75{\cdot}r_1$ signed \\
\bottomrule
\end{tabular}

\textbf{Rule:} LSR $\Rightarrow$ unsigned only. ASR $\Rightarrow$ signed. LSL $\Rightarrow$ both. Decompose $k$ as $1\pm2^{-n}$ or $2^n\pm1$ to fit one ADD/SUB/RSB.

\begin{qp}\textbf{Q3 P4} --- $A=k{\cdot}B/8$ (signed, one instr):
\begin{itemize}
\item $6B/8$: \asm{sub r0,r1,r1,asr \#2}  ($B - B/4$)
\item $7B/8$: \asm{sub r0,r1,r1,asr \#3}
\item $9B/8$: \asm{add r0,r1,r1,asr \#3}
\item $10B/8$: \asm{add r0,r1,r1,asr \#2}
\end{itemize}
\end{qp}

\subsection{S-suffix shifts \& RRX}
S-variants of LSL/LSR/ASR/ROR update C = \textbf{last bit shifted out}; also N,Z.

\asm{RRX Rd,Rn}: 33-bit rotate right by 1; C$\to$b31, b0$\to$(discarded unless RRXS). RRXS takes b0$\to$C.

\textbf{C extraction trick:} \asm{lsls r0,\#1} puts old b31 into C. \asm{lsrs r0,\#n} puts old b($n-1$) into C.

% ==========================================================
\section{Bitwise Logic (L16--17)}
\begin{tabular}{@{}ll@{}}
\toprule
Insn & Effect \\
\midrule
\asm{AND Rd,Rn,Op2} & $Rn \mathop{\&} Op2$ \\
\asm{ORR Rd,Rn,Op2} & $Rn \mathop{|} Op2$ (set bits) \\
\asm{EOR Rd,Rn,Op2} & $Rn \oplus Op2$ (toggle) \\
\asm{BIC Rd,Rn,Op2} & $Rn \mathop{\&} \sim Op2$ (clear) \\
\asm{ORN Rd,Rn,Op2} & $Rn \mathop{|} \sim Op2$ \\
\asm{MVN Rd,Op2}    & $\sim Op2$ \\
\asm{BFC Rd,\#s,\#w}& clear $w$ bits at bit $s$ \\
\asm{BFI Rd,Rn,\#s,\#w} & insert $Rn[w{-}1{:}0]$ into $Rd[s{+}w{-}1{:}s]$ \\
\bottomrule
\end{tabular}

\textbf{BIC has no C-operator equivalent} --- must write \cee{a \& \textasciitilde mask}.

\subsection{Bit-field idioms}
\textbf{Replace a nibble/field} (``clear-then-insert''):
\begin{lstlisting}
@ set bits[s+w-1:s] = v (w-bit)
bic r0,r0,mask,lsl #s  @ clear
orr r0,r0,val, lsl #s  @ insert
\end{lstlisting}
C form: \cee{a \&= \textasciitilde(mask{\mbox{<}}{\mbox{<}}s); a |= (val \& mask){\mbox{<}}{\mbox{<}}s;}

\textbf{Insert digit at an end} (after shifting): the vacated nibble is already zero, so no BIC needed --- just ORR.

\begin{hwp}\textbf{HW7 P4:} replace $H_1$ with $0xA$ in \cee{uint16\_t a}:
\begin{lstlisting}
ldr r1,=15
bic r0,r0,r1,lsl #4
ldr r2,=10
orr r0,r0,r2,lsl #4
\end{lstlisting}
\end{hwp}

\begin{hwp}\textbf{HW7 P5} (hand-eval): R0=0xAD, R1=0x03. \asm{bic r3,r0,r1}=0xAC; \asm{orr r4,r3,r1}=0xAF; \asm{eor r5,r0,r1}=0xAE.\end{hwp}

\begin{hwp}\textbf{HW7 P6b} (shift-left + insert low digit, $H_3H_2H_1H_0\to H_2H_1H_0\,3$):
\begin{lstlisting}
lsl r0,#4
orr r0,#3
\end{lstlisting}
\textbf{P6c} (shift-right + insert top digit, $\to\,7\,H_3H_2H_1$):
\begin{lstlisting}
lsr r0,#4
ldr r1,=7
lsl r1,#12
orr r0,r1
\end{lstlisting}
\end{hwp}

\begin{qp}\textbf{Q3 P5} --- ASM uses \asm{bic...lsl \#8}+\asm{orr...lsl \#8}, mask=15: replace $H_2$ (bits 11:8) with low nibble of $b$.
\cee{a \&= \textasciitilde(15{\mbox{<}}{\mbox{<}}8); a |= (b \& 15){\mbox{<}}{\mbox{<}}8;}\end{qp}

% ==========================================================
\section{Flags NZCV (L17--18)}
APSR: N=b31, Z=b30, C=b29, V=b28 (Q=b27).\\
\textbf{N,Z} used by both signed and unsigned. \textbf{C} = unsigned carry/borrow. \textbf{V} = signed overflow.

Processor sets ALL flags every arithmetic --- it doesn't know signed vs unsigned. Programmer picks CCS.

\subsection{Hand computation (N-bit system)}
Let $C_N=2^N$, $U_{\max}=2^N{-}1$, $S_{\min}=-2^{N-1}$, $S_{\max}=2^{N-1}{-}1$.

\textbf{Per operation:}
\begin{enumerate}
\item \textbf{Unsigned} view: map negatives via $+C_N$, compute, reduce mod $C_N$. Cross-boundary $\Rightarrow$ \textbf{C=1} (add) or \textbf{C=0} (sub borrow).
\item \textbf{Signed} view: compute directly. If outside $[S_{\min},S_{\max}]$ $\Rightarrow$ \textbf{V=1}, correct by $\pm C_N$.
\item \textbf{Binary} result is the SAME for both views. \textbf{N}=MSB. \textbf{Z}=(result==0).
\end{enumerate}

\begin{hwp}\textbf{HW8 P1} (5-bit, $C_N=32$) ${-}15+({-}5)$: U: $17{+}27{=}44{-}32{=}12$, \textbf{C=1}. S: ${-}20{<}{-}16$, $+32{=}12$, \textbf{V=1}. Binary 0b01100 (N=0, Z=0).\end{hwp}

\begin{qp}\textbf{Q4 P2} (5-bit) $16{-}18$: U: $16{-}18{=}{-}2{+}32{=}30$, borrow so \textbf{C=0}. S: $16{-}({-}14){=}30{>}15$ wait -- actually $16$ exceeds $S_{\max}$ itself so... the exam uses $16{-}18$ with signed view: ${-}16{-}({-}14){=}{-}2\in[{-}16,15]$, \textbf{V=0}. Binary 0b11110 (N=1, Z=0).\end{qp}

\textbf{Mnemonic ``C=!Borrow'':} after \asm{subs}: no borrow $\to$ C=1; borrow $\to$ C=0.

\subsection{CMP / CMN / TST / TEQ}
\begin{tabular}{@{}lll@{}}
\toprule
Insn & Op (discarded) & Flags \\
\midrule
CMP Rn,Op2 & $Rn-Op2$ & N,Z,C,V \\
CMN Rn,Op2 & $Rn+Op2$ & N,Z,C,V \\
TST Rn,Op2 & $Rn\mathop{\&}Op2$ & N,Z,C \\
TEQ Rn,Op2 & $Rn\oplus Op2$ & N,Z,C \\
\bottomrule
\end{tabular}

TST/TEQ do NOT affect V. CMP $\equiv$ SUBS-no-writeback; CMN $\equiv$ ADDS-no-writeback.

\textbf{Pseudo:} \asm{cmp r0,\#-5} assembles as \asm{CMN r0,\#5}. Signed CCS still valid.

\textbf{APSR access:} \asm{mrs r0,apsr} (read); \asm{msr apsr\_nzcvq,r0} (write).

% ==========================================================
\section{CCS Table (L18--19)}
\begin{tabular}{@{}lccl@{}}
\toprule
CCS & S/U & Flags \\
\midrule
EQ / NE  & -- & $Z{=}1$ / $Z{=}0$ \\
GE       & S  & $V{=}N$ \\
GT       & S  & $V{=}N \wedge Z{=}0$ \\
LT       & S  & $V{\ne}N$ \\
LE       & S  & $V{\ne}N \vee Z{=}1$ \\
HS (CS)  & U  & $C{=}1$ \\
LO (CC)  & U  & $C{=}0$ \\
HI       & U  & $C{=}1 \wedge Z{=}0$ \\
LS       & U  & $C{=}0 \vee Z{=}1$ \\
\bottomrule
\end{tabular}

\textbf{C$\to$CCS mapping:}
\begin{tabular}{@{}lcc@{}}
\toprule
C & Signed & Unsigned \\
\midrule
\cee{==} & EQ & EQ \\
\cee{!=} & NE & NE \\
\cee{>}  & GT & HI \\
\cee{>=} & GE & HS \\
\cee{<}  & LT & LO \\
\cee{<=} & LE & LS \\
\bottomrule
\end{tabular}

\textbf{Mnemonic ``HiLow unSun'':} Higher/Lower/Same = unsigned (can't see the Sun above/below horizon).

\textbf{Consistency rule:} C-type determines \textbf{CCS family} (GT/GE/LT/LE vs HI/HS/LO/LS), \textbf{shift} (ASR vs LSR), and \textbf{load suffix} (SB/SH vs B/H).

\begin{hwp}\textbf{HW8 P2:} CMP with equal-magnitude diffs give identical flags because $(-20)+64-((-22)+64)=22-20$. CCS is \textbf{flag-driven}, not value-driven.\end{hwp}

\begin{hwp}\textbf{HW8 P3} (max via NL): $r_0 \ge r_1$?
\begin{lstlisting}
cmp r0,r1
blt .else        @ !(a>=b)
bx  lr           @ then: r0 is a
.else:
mov r0,r1
bx  lr
\end{lstlisting}
\end{hwp}

% ==========================================================
% ============  PART II: HW9-12 / QUIZ 4 MATERIAL  =========
% ==========================================================

\section{LDR/STR Matrix (L20--21)}
Two regs: \asm{Ra}=addr, \asm{Rv}=data.

\subsection{LOAD suffixes}
\begin{tabular}{@{}llll@{}}
\toprule
Suf & Size & Ext & C type \\
\midrule
---  & word  & ---  & \cee{(u)int32\_t} \\
B    & ubyte & zero & \cee{uint8\_t} \\
SB   & sbyte & sign & \cee{int8\_t} \\
H    & uhalf & zero & \cee{uint16\_t} \\
SH   & shalf & sign & \cee{int16\_t} \\
\bottomrule
\end{tabular}

\textbf{STORE has NO sign variant} --- \asm{STR/STRH/STRB} only.

\subsection{Three addressing modes}
\begin{tabular}{@{}ll@{}}
\toprule
Mode & Effect \\
\midrule
\asm{[Ra,\#os]}   & load from $Ra{+}os$; Ra unchanged \\
\asm{[Ra,\#os]!}  & $Ra{+}{=}os$, then load \\
\asm{[Ra],\#os}   & load, then $Ra{+}{=}os$ \\
\bottomrule
\end{tabular}

\textit{Mnemonic:} \texttt{!}=BEFORE; trailing \texttt{\#os}=AFTER.

\subsection{Register-offset (L21)}
\begin{lstlisting}
ldr  Rv,[Ra,Ri,LSL #2]  @ word, Ra+4*Ri
ldrh Rv,[Ra,Ri,LSL #1]  @ half, Ra+2*Ri
ldrb Rv,[Ra,Ri]          @ byte, Ra+Ri
\end{lstlisting}
Shift $=\log_2$(elem size). No pre/post writeback form.

\subsection{LDRD / STRD --- 64-bit (L23)}
\begin{lstlisting}
ldrd Rv_L,Rv_H,[Ra,#os]   @ basic
ldrd Rv_L,Rv_H,[Ra,#os]!  @ pre
ldrd Rv_L,Rv_H,[Ra],#os   @ post
\end{lstlisting}
L=LSW (low addr), H=MSW, little-endian. \textbf{No reg-offset form.}

\begin{hwp}\textbf{HW9 P1} (16 bytes 0x00..0x0F @ 0x20000000):\\
\asm{ldr r4,=0x20000000} $\to$ r4=base.\\
\asm{ldr r0,[r4,\#4]}: basic $\to$ r0=0x07060504.\\
\asm{ldr r1,[r4,\#2]!}: pre $\to$ r4=0x20000002, r1=0x05040302.\\
\asm{ldr r2,[r4],\#4}: post $\to$ r2=0x05040302, then r4=0x20000006.\\
\asm{ldr r3,[r4]}: r3=0x09080706.\end{hwp}

\begin{hwp}\textbf{HW10 P2 T1:} \asm{ldrd r2,r3,[r0,\#8]!} on \cee{ipt[i]=i<{<}4} $\to$ r0+=8; r2=0xB0A09080, r3=0xF0E0D0C0.\end{hwp}

% ==========================================================
\section{C Pointer $\to$ ASM (L20, 22)}

\subsection{32-bit ptr: step=4}
\begin{tabular}{@{}lll@{}}
\toprule
C & LDR & STR \\
\midrule
\cee{*(p+k)} & \asm{[r1,\#4k]}  & \asm{[r1,\#4k]} \\
\cee{*++p}   & \asm{[r1,\#4]!}  & \asm{[r1,\#4]!} \\
\cee{*p++}   & \asm{[r1],\#4}   & \asm{[r1],\#4}  \\
\cee{*--p}   & \asm{[r1,\#-4]!} & \asm{[r1,\#-4]!}\\
\cee{*p-{}-} & \asm{[r1],\#-4}  & \asm{[r1],\#-4} \\
\bottomrule
\end{tabular}

\subsection{16-bit signed: step=2 $\to$ LDRSH}
\cee{*(p+k)} $\to$ \asm{ldrsh r0,[r1,\#2k]}\\
\cee{*++p}   $\to$ \asm{ldrsh r0,[r1,\#2]!}\\
\cee{*p++}   $\to$ \asm{ldrsh r0,[r1],\#2}

Unsigned: LDRSH$\to$LDRH. Stores: STRH (no sign).

\subsection{8-bit: step=1 $\to$ LDRB/LDRSB}
\cee{*(p+k)} $\to$ \asm{ldrb r0,[r1,\#k]}\\
\cee{*p++}   $\to$ \asm{ldrb r0,[r1],\#1}\\
\cee{*++p}   $\to$ \asm{ldrb r0,[r1,\#1]!}

\subsection{Array indexing (const ptr, index i)}
\begin{lstlisting}
ldr  r4,[r0,r3,lsl #2]  @ 32b: pArr[i]
ldrh r4,[r0,r3,lsl #1]  @ 16b
ldrb r4,[r0,r3]          @  8b
\end{lstlisting}

\begin{hwp}\textbf{HW9 P2 mp\_task1:} \asm{ldrsh r2,[r0],\#4} then \asm{str r2,[r1,\#0]}: \cee{*p32 = (int32\_t)*p16s; p16s+=2;} (post-inc 4 = skip halfword).\end{hwp}

\begin{hwp}\textbf{HW9 P2 mp\_task2:} \asm{ldrh r2,[r0,\#2]!}: \cee{p16u++; *p32 = (int32\_t)*p16u;}.\end{hwp}

\begin{hwp}\textbf{HW10 P2 T4:} \asm{ldrh r3,[r0,r2,LSL \#1]} with r2=1,2,3 $\to$ iptr[1]=0x3020, iptr[2]=0x5040. Register-offset = typed array indexing.\end{hwp}

\subsection{HW9 P3 four-task register state}

Entry: R0=\texttt{ipt} (0x20001010), R1=\texttt{opt}. Words: +0=0x30201000, +4=0x70605040, +8=0xB0A09080, +12=0xF0E0D0C0.

\begin{tabular}{@{}llll@{}}
\toprule
Task & Form & R0 after & Printout \\
\midrule
T1 & basic \asm{[r0,\#4]} \& \asm{\#0xC} & unchanged & 0x70605040 / 0xF0E0D0C0 \\
T2 & pre-idx \asm{\#4!} $\times 2$ & 0x...1018 & 0x70605040 / 0xB0A09080 \\
T3 & post-idx \asm{\#4} $\times 2$ & 0x...1018 & 0x30201000 / 0x70605040 \\
T4 & reg-off \asm{[r0,r2,LSL \#2]} & unchanged & 0x70605040 / 0xB0A09080 \\
\bottomrule
\end{tabular}

C equiv (all use \cee{int32\_t *p=(int32\_t*)ipt}):
\begin{itemize}
\item T1: \cee{opt[0]=p[1]; opt[1]=p[3];}
\item T2: \cee{p++; opt[0]=*p; p++; opt[1]=*p;}
\item T3: \cee{opt[0]=*p; p++; opt[1]=*p; p++;}
\item T4: \cee{int i=1; opt[0]=p[i]; i++; opt[1]=p[i]; i++;}
\end{itemize}

\subsection{HW10 P1 four-line memory trace}

\cee{uint8\_t arr[i]=2*i} @ 0x20000100. \texttt{p8a}=0x100, \texttt{p8b}=0x10C.

\begin{tabular}{@{}lllll@{}}
\toprule
Ln & C RHS & Mem $\Delta$ & p8a & p8b \\
\midrule
4 & \cee{*(p8a+2)} & --- & 0x100 & 0x10D \\
5 & \cee{++*(p8a+4)} & [0x104]: 0x08$\to$0x09 & 0x100 & 0x10E \\
6 & \cee{++*p8a++} & [0x100]: 0x00$\to$0x01 & 0x101 & 0x10F \\
7 & \cee{*p8a++} & --- & 0x102 & 0x110 \\
\bottomrule
\end{tabular}

Final 4 bytes at 0x20000{\textunderscore}10C: 0x04, 0x09, 0x01, 0x02.

\subsection{HW10 P2 four-task LDRD/LDRSB/LDRSH trace}

\cee{ipt[i]=i<{<}4} @ 0x20001010. Entry R0=0x...1010.

\begin{tabular}{@{}lll@{}}
\toprule
Task & ASM (1 key line) & Result \\
\midrule
T1 & \asm{ldrd r2,r3,[r0,\#8]!} & R0=0x...1018, r2=0xB0A09080, r3=0xF0E0D0C0 \\
T2 & \asm{ldrsb r?,[r0,\#4]!} $\times 2$ & r2=0x00000040; r3=\textbf{0xFFFFFF80} \\
T3 & \asm{ldrsh ...,\#10/\#4} post & r2=0x00001000; r3=\textbf{0xFFFFB0A0}. R0=0x...101E (unaligned!) \\
T4 & \asm{ldrh r3,[r0,r2,LSL \#1]} & r2=1$\to$3; r3=0x5040 \\
\bottomrule
\end{tabular}

Sign-extension pivots: byte 0x80 (b7=1); half 0x8000/0xB0A0 (b15=1). Byte 0x40: b7=0 $\to$ zero-like extension.

% ==========================================================
\section{Sign-Extension Pitfalls}
\begin{tabular}{@{}lll@{}}
\toprule
Insn & 0x80 $\to$ & 0xB0A0 $\to$ \\
\midrule
LDRB  & 0x00000080          & --- \\
LDRSB & \textbf{0xFFFFFF80} & --- \\
LDRH  & ---                 & 0x0000B0A0 \\
LDRSH & ---                 & \textbf{0xFFFFB0A0} \\
\bottomrule
\end{tabular}

Check MSBit of loaded byte/half to predict extension.

\begin{hwp}\textbf{HW10 P2 T2:} \asm{ldrsb r2,[r0,\#4]!} at byte 0x40 $\to$ r2=0x00000040; next \asm{ldrsb} at 0x80 $\to$ r3=\textbf{0xFFFFFF80}.\end{hwp}

\begin{hwp}\textbf{HW10 P2 T3:} \asm{ldrsh r3,[r0],\#4} at halfword 0xB0A0 $\to$ r3=\textbf{0xFFFFB0A0}.\end{hwp}

\begin{qp}\textbf{Q3 P6} (EABI sign-ext): \cee{int8\_t 0x81} as fn arg $\to$ R1 = 0xFFFFFF81. \cee{int8\_t 0x7E} $\to$ R0 = 0x0000007E.\end{qp}

% ==========================================================
\section{Inc/Dec --- Values vs Ptrs (L22)}
\textbf{C precedence:} L1 (LTR): postfix \cee{++/--}, \cee{[]}, \cee{()}. L2 (RTL): prefix \cee{++/--}, \cee{*}, \cee{\&}, cast, \cee{sizeof}.

\cee{p++} binds tighter than adjacent \cee{*} or prefix \cee{++}, but its side effect (increment) happens \textbf{after} the original p is used.

\subsection{Prefix vs postfix on a value}
\cee{B[0] = ++A[0];} $\to$ inc first, then assign:
\begin{lstlisting}
ldr r2,[r0] / add r2,#1
str r2,[r0] / str r2,[r1]
\end{lstlisting}

\cee{B[1] = A[1]++;} $\to$ assign old, then inc:
\begin{lstlisting}
ldr r2,[r0,#4] / str r2,[r1,#4]
add r2,#1       / str r2,[r0,#4]
\end{lstlisting}

\subsection{Composite patterns}
\begin{tabular}{@{}ll@{}}
\toprule
C expr & ASM \\
\midrule
\cee{*p++}    & \asm{ldrb r2,[r0],\#1} \\
\cee{*++p}    & \asm{ldrb r2,[r0,\#1]!} \\
\cee{*p++=v}  & \asm{strb r2,[r0],\#1} \\
\cee{++*p}    & ldrb / add\#1 / strb \\
\cee{++*p++}  & ldrb[r0],\#1 / add / strb[r0,\#-1] \\
\cee{*p=*q++} & ldrb [r1],\#1 / strb [r0] \\
\bottomrule
\end{tabular}

\begin{hwp}\textbf{HW10 P1 L5:} \cee{*p8b++ = ++*(p8a+4);} --- inc byte at p8a+4 (0x08$\to$0x09), write back, store:
\begin{lstlisting}
ldrb r2,[r0,#4]
add  r2,r2,#1
strb r2,[r0,#4]    @ write-back!
strb r2,[r1],#1
\end{lstlisting}\end{hwp}

\begin{hwp}\textbf{HW10 P1 L6:} \cee{*p8b++ = ++*p8a++;}:
\begin{lstlisting}
ldrb r2,[r0],#1    @ r0 now at old+1
add  r2,r2,#1
strb r2,[r0,#-1]   @ write back old addr
strb r2,[r1],#1
\end{lstlisting}\end{hwp}

\textbf{Ptr arithmetic:} \cee{p8++}$\to+1$, \cee{p16++}$\to+2$, \cee{p32++}$\to+4$. \cee{p-arr} typed $\to$ elements; \cee{(void*)p-(void*)arr} $\to$ bytes.

% ==========================================================
\section{Function Calls (L23)}
\textbf{AAPCS regs:}
\begin{itemize}
\item \textbf{R0--R3}: args 1--4, ret (R0 or R1:R0); caller-saved.
\item \textbf{R12 (IP)}: scratch; caller-saved.
\item \textbf{R4--R11}: general; \textbf{callee-saved} (push/pop if used).
\item \textbf{LR}: ret addr; save only if stem.
\end{itemize}

\textbf{Leaf} (no \asm{bl}): \asm{bx lr}; no LR push.\\
\textbf{Stem}: \asm{push \{lr\}} + \asm{pop \{pc\}}.

\textbf{BL:} \asm{bl tgt} $\Rightarrow$ LR=PC\_next; PC=tgt.

\subsection{Stem template}
\begin{lstlisting}
my_stem:
  push {r4-r7,lr}
  mov  r4,r0     @ stash args
  mov  r5,r1
  mov  r6,r2
  mov  r7,r3
  @ prep R0-R3, bl inner
  pop  {r4-r7,pc}
\end{lstlisting}

\subsection{Load args by C type}
\begin{tabular}{@{}ll@{}}
\toprule
C type & Load \\
\midrule
\cee{(u)int32\_t / *} & \asm{ldr} \\
\cee{int16\_t}        & \asm{ldrsh} \\
\cee{uint16\_t}       & \asm{ldrh} \\
\cee{int8\_t/char}    & \asm{ldrsb} \\
\cee{uint8\_t}        & \asm{ldrb} \\
\cee{(u)int64\_t}     & \asm{ldrd rL,rH,...} \\
\bottomrule
\end{tabular}

\textbf{Return:} $\le 32b$ $\to$ STR/STRH/STRB on R0. 64b $\to$ \asm{strd r0,r1,...}.

\subsection{Array-loop + call}
\begin{lstlisting}
  ldr r8,=0
lp:
  cmp r8,r7
  bge end
  ldr r0,[r4,r8,lsl #2]  @ x[i]
  ldr r1,[r5,r8,lsl #2]  @ y[i]
  bl  mp_umod_32bit_s
  str r0,[r6,r8,lsl #2]  @ z[i]
  add r8,#1
  b   lp
end:
\end{lstlisting}

64b elts: \asm{ldrd r0,r1,[r4],\#8}/\asm{ldrd r2,r3,[r5],\#8}/\asm{bl}/\asm{strd r0,r1,[r6],\#8}.

% ==========================================================
\section{IT Blocks \& CBZ/CBNZ (L24)}
\asm{IT\{x\{y\{z\}\}\} cond} --- up to 4 slots. First=T. Later T=cond, E=inverse. Forms: IT, ITT, ITE, ITTT, ITTE, ITEE, ITTEE, ITTTT.

\subsection{Hard rules}
\begin{enumerate}
\item Max \textbf{4} cond instrs.
\item First slot = cond; T=cond, E=inverse.
\item \textbf{Cond branch (Bcc) only in LAST slot.}
\item \textbf{CBZ/CBNZ FORBIDDEN inside IT.}
\item CBZ/CBNZ: forward only ($\le$126 B), Rn = R0--R7, no flags.
\end{enumerate}

\subsection{CBZ/CBNZ}
\begin{lstlisting}
cbz  Rn,label    @ if Rn==0
cbnz Rn,label    @ if Rn!=0
\end{lstlisting}
Replaces \asm{cmp Rn,\#0/beq}. Perfect for null-terminator after \asm{ldrb}.

\subsection{CEX if/else (L25)}
\cee{if (x{\mbox{>}}=0) x{\mbox{<}}{\mbox{<}}=1; else x{\mbox{>}}{\mbox{>}}=4;}
\begin{lstlisting}
cmp   r0,#0
ite   ge
lslge r0,#1
asrlt r0,#4
bx    lr
\end{lstlisting}

\subsection{CEX, 2 instrs/branch}
\begin{lstlisting}
  cmp     r0,r1
  ittee   le       @ T,T=le ; E,E=gt
  ldrle   r0,[r3],#4
  asrle   r0,r0,#1
  ldrshgt r0,[r2,#2]!
  lslgt   r0,r0,#1
  bx      lr
\end{lstlisting}
\begin{hwp}\textbf{HW11 P1 P3:} exact ITTEE le pattern above.\end{hwp}

\subsection{Cond-B to short-circuit IT}
\begin{lstlisting}
cmp   r0,#-5     @ = CMN r0,#5
itt   lt
ldrlt r1,=-10    @ T1
blt   end_label  @ T2: LAST slot
\end{lstlisting}

\begin{hwp}\textbf{HW11 P2 P3} (piecewise $f(x)$):
\begin{lstlisting}
  cmp   r0,#-5
  itt   lt
  ldrlt r1,=-10
  blt   end
  cmp   r0,#5
  ite   ge
  ldrge r1,=10
  lsllt r1,r0,#1
end:
  mov   r0,r1
  bx    lr
\end{lstlisting}\end{hwp}

% ==========================================================
\section{If/If-Else (L25)}

\subsection{Positive Logic (PL)}
Branch-to-then on C TRUE. \textbf{Same} CCS.
\begin{lstlisting}
  cmp r0,r1
  bgt pl_then
pl_else:
  @ else
  b  pl_end
pl_then:
  @ then
pl_end:
  bx lr
\end{lstlisting}

\subsection{Negative Logic (NL)}
Branch-to-else on C FALSE. \textbf{Opposite} CCS. NL reads closer to C.
\begin{lstlisting}
  cmp r0,r1
  ble nl_else    @ LE = !GT
nl_then:
  @ then
  b  nl_end
nl_else:
  @ else
nl_end:
  bx lr
\end{lstlisting}

\begin{hwp}\textbf{HW11 P1} same C in 3 styles: \cee{if (x{\mbox{>}}y) return *(++ptrA)*2; else return *(ptrB++)/2;}
--- Part 1 PL: \asm{cmp/bgt pl\_then}; Part 2 NL: \asm{cmp/ble nl\_else}; Part 3 CEX: ITTEE le. \cee{int16\_t *} $\to$ LDRSH, signed $/2$ $\to$ ASR.\end{hwp}

\subsection{If-elseif-else (CMP chain)}
\begin{lstlisting}
  cmp r0,#-5
  bge elseif     @ NL: x<-5
then:
  ldr r0,=-10
  b   end
elseif:
  cmp r0,#5
  blt else_blk   @ NL: x>=5
  ldr r0,=10
  b   end
else_blk:
  lsl r0,r0,#1   @ 2*x
end:
  bx  lr
\end{lstlisting}
\begin{hwp}\textbf{HW11 P2 P2:} exactly this 3-way cascade.\end{hwp}

\subsection{Compound OR}
\cee{if (x{\mbox{<}}=20 || x{\mbox{>}}=25) then else else}
\begin{lstlisting}
cmp r0,#20
ble then_block
cmp r0,#25
bge then_block
b   else_block
\end{lstlisting}
\begin{hwp}\textbf{HW11 P3 P1:} literal of above. \cee{*++ptrA}$\to$LDRH pre-inc, $*4$$\to$LSL\#2; \cee{*++ptrB}$\to$LDR pre-inc, unsigned $/4$$\to$LSR\#2.\end{hwp}

\subsection{Compound AND (De Morgan)}
\cee{if (x{\mbox{>}}0 \&\& x{\mbox{<}}8) then else else}
\begin{lstlisting}
cmp r0,#0
ble else_block   @ x<=0 -> else
cmp r0,#8
bge else_block   @ x>=8 -> else
@ fall through = then
\end{lstlisting}

\begin{hwp}\textbf{HW11 P3 P2 CEX:} first sub-test can't fit in IT; \asm{bgt check25}, then ITTEE ge:
\begin{lstlisting}
  cmp   r0,#20
  bgt   check25
  ldrh  r0,[r1,#2]!
  lsl   r0,r0,#2
  bx    lr
check25:
  cmp   r0,#25
  ittee ge
  ldrhge r0,[r1,#2]!
  lslge  r0,r0,#2
  ldrlt  r0,[r2,#4]!
  lsrlt  r0,r0,#2
  bx     lr
\end{lstlisting}\end{hwp}

% ==========================================================
\section{Loops (L26)}
Labels: \texttt{\_lp}=top, \texttt{\_end}=after.

\subsection{while (top-tested)}
\begin{lstlisting}
  ldrb r2,[r1],#1   @ PRIMING
lp:
  cbz  r2,end
  @ body
  ldrb r2,[r1],#1   @ in-loop (DUP)
  b    lp
end:
\end{lstlisting}

\subsection{do-while (bottom-tested)}
\begin{lstlisting}
lp:
  @ body
  cmp r0,r1
  blt lp
\end{lstlisting}

\subsection{for loop}
\begin{lstlisting}
  mov r4,#0
lp:
  cmp r4,r5
  bge end
  @ body
  add r4,#1
  b   lp
end:
\end{lstlisting}

\subsection{while(1) + break (no priming dup)}
\begin{lstlisting}
lp:
  ldrb r2,[r1],#1   @ SINGLE read
  cbz  r2,end
  @ body
  b    lp
end:
\end{lstlisting}

\begin{hwp}\textbf{HW12 P2 P1} (plain while, lc$\to$UC copy): priming LDRB + in-loop LDRB with \asm{sub r2,\#0x20} between; terminator stored at \texttt{\_end}.\end{hwp}

\begin{hwp}\textbf{HW12 P2 P3} (while(1)+break): single in-loop LDRB:
\begin{lstlisting}
lp:
  ldrb r2,[r1],#1
  cbz  r2,end
  sub  r2,r2,#0x20
  strb r2,[r0],#1
  b    lp
end:
  strb r2,[r0],#1  @ '\0' store
  bx   lr
\end{lstlisting}\end{hwp}

\begin{hwp}\textbf{HW12 P1} (swap-last-2-of-3): r12=pair counter (scratch, no push):
\begin{lstlisting}
  mov  r12,#0
lp:
  ldrb r1,[r0],#1
  cbz  r1,end
  ldrb r2,[r0],#1
  cbz  r2,end
  ldrb r3,[r0],#1
  cbz  r3,end
  strb r3,[r0,#-2] @ *(s-2)=c3
  strb r2,[r0,#-1] @ *(s-1)=c2
  add  r12,r12,#1
  b    lp
end:
  mov  r0,r12
  bx   lr
\end{lstlisting}
3 CBZ after 3 post-index LDRB catches early null.\end{hwp}

\subsection{while-using-CEX (rare)}
\begin{lstlisting}
lp:
  cmp  r0,r1
  it   lt
  blt  body_in_it  @ last slot
  b    end
body_in_it:
  @ body
  b lp
end:
\end{lstlisting}

% ==========================================================
\section{break, continue, switch (L27)}
\begin{itemize}
\item \cee{break} = forward \asm{b end} (or inline CBZ/CBNZ).
\item \cee{continue} (while/do-while) = back \asm{b lp}.
\item \cee{continue} in \texttt{for} = branch to \textbf{update}, not top.
\item CBZ/CBNZ: no IT; free-standing preferred null-test.
\end{itemize}

\subsection{switch: CMP/BEQ ladder}
\begin{lstlisting}
  cmp r0,#65      @ 'A'
  beq case_a
  cmp r0,#66
  beq case_b
  b   default_lbl
case_a:
case_b:           @ stacked=fall-through
  ldr r0,=80
  b   end
default_lbl:
  ldr r0,=69
end:
  bx lr
\end{lstlisting}

\subsection{switch: ADR+TBB (O(1))}
\begin{lstlisting}
  adr r1,table
  tbb [r1,r0]     @ PC=PC+4+2*table[r0]
case_a:
  ...
table:
  .byte 0
  .byte (case_b-case_a)/2
\end{lstlisting}

% ==========================================================
% =========  PART III: PAST QUIZZES PROBLEM-BANK  =========
% ==========================================================

\section{Quiz 3 Problems (ASM in C)}

\begin{qp}\textbf{Q3 P1} Return signed -2 from int32 fn: return reg R0 = \textbf{0xFFFFFFFE} (two's comp). Not 0x80000002.\end{qp}

\begin{qp}\textbf{Q3 P2} 64-bit add test design: pick A, B whose \textbf{low 32b} sum crosses $2^{32}$ for one case and doesn't for the other. Shifting a nibble into low-word MSB flips C.
A1=0x377{\mbox{<}}{\mbox{<}}24, B1=0x287{\mbox{<}}{\mbox{<}}24 $\to$ C=0.
A2=0x378{\mbox{<}}{\mbox{<}}24, B2=0x288{\mbox{<}}{\mbox{<}}24 $\to$ C=1.\end{qp}

\begin{qp}\textbf{Q3 P3} \asm{push \{r4,r8,r6\}}: SP$-$=12; R4@..14, R6@..18, R8@..1C. Byte@..19 = 2nd byte of R6 little-endian.\end{qp}

\begin{qp}\textbf{Q3 P4} $A=k{\cdot}B/8$ single-instr w/ ASR Op2 (see Barrel Shifter sect).\end{qp}

\begin{qp}\textbf{Q3 P5} BIC/ORR bitfield splice $\leftrightarrow$ C: ASM \asm{bic r0,r0,\#15,lsl \#8} + \asm{orr r0,r0,(b\&15),lsl \#8} = \cee{a \&= \textasciitilde(15{\mbox{<}}{\mbox{<}}8); a |= (b\&15){\mbox{<}}{\mbox{<}}8;}\end{qp}

\begin{qp}\textbf{Q3 P6} sign-ext at call: \cee{int8\_t 0x7E}$\to$R0=0x7E; \cee{int8\_t 0x81}$\to$R1=0xFFFFFF81.\end{qp}

% ==========================================================
% [SEMESTER QZ4] Spring-2026 Quiz 4 (qz-26a) - LDR/STR + flow ctrl
% Source: qz-26a-qz4-ldr-str-n-flow-ctrl-soln-26-04.pdf
\section{Quiz 4 (Spring 2026, 105 pts)}

Memory at 0x2000\_0000: bytes \texttt{88 89 8A 8B 8C 8D 8E 8F 70 71 72 73 74 75 76 77}.

\begin{qp}\textbf{P1a (8pts)} task1: \asm{ldrsh r2,[r0,\#4]} loads halfword 0x8D8C $\to$ r2 = \textbf{0xFFFF8D8C} (sign-ext, MSB=1). \asm{ldrsh r3,[r0,\#8]} loads 0x7170 $\to$ r3 = 0x00007170.\end{qp}

\begin{qp}\textbf{P1b (8pts)} task2: \asm{ldrsb r2,[r0,\#2]!} byte 0x8A $\to$ r2 = \textbf{0xFFFFFF8A}, r0=0x...0002. \asm{ldrsb r3,[r0,\#4]!} byte at 0x...0006 = 0x8E $\to$ r3 = \textbf{0xFFFFFF8E}, r0=0x...0006.\end{qp}

\begin{qp}\textbf{P1c (8pts)} After task1 stores (r2 then r3 with post-idx \#4) at 0x2000\_0020: bytes \texttt{8C 8D FF FF | 70 71 00 00}.\end{qp}

\begin{qp}\textbf{P1d (6pts)} \cee{val = *(ptr+i)} for \cee{uint16\_t*}:
\begin{lstlisting}
ldrh r2,[r0,r1,lsl #1]   @ scale i*2
\end{lstlisting}
LSL=$\log_2$(sizeof). Index in elements; addr in bytes.\end{qp}

\begin{qp}\textbf{P2 (30pts)} Three C-pointer ops on \cee{int16\_t *src,*dst}:
\begin{lstlisting}
@ *dst++ = *src++;
ldrsh r2,[r0],#2
strh  r2,[r1],#2
@ *dst++ = *++src;
ldrsh r2,[r0,#2]!
strh  r2,[r1],#2
@ *dst++ = ++*src;     @ inc memory!
ldrsh r2,[r0]
add   r2,#1
strh  r2,[r0]          @ writeback
strh  r2,[r1],#2
\end{lstlisting}
\textbf{Trap:} \cee{++*src} modifies the dereferenced value at memory --- needs explicit STRH writeback.\end{qp}

\begin{qp}\textbf{P3 (20pts)} LDRD/STRD round trip. \cee{ipt[i]=i<{<}4} $\to$ words 0x30201000, 0x70605040, 0xB0A09080, 0xF0E0D0C0.
\begin{lstlisting}
ldrd  r2,r3,[r0,#8]!   @ r2=0xB0A09080, r3=0xF0E0D0C0
str   r3,[r1,#0]        @ opt[0]=0xF0E0D0C0
str   r2,[r1,#4]        @ opt[1]=0xB0A09080
sub   r0,r0,#8          @ rewind
ldrsh r3,[r0],#4        @ r3=0x00001000
ldrsh r2,[r0,#4]!       @ r2=0xFFFF9080 (sign-ext)
strd  r2,r3,[r1,#8]!    @ opt[2..3]
bx    lr
\end{lstlisting}
Outputs: 0xF0E0D0C0 / 0xB0A09080 / 0xFFFF9080 / 0x1000.\\
\textbf{Trap:} LDRD Rt=low addr (LSW), Rt2=high addr (MSW). Order matters.\end{qp}

\begin{qp}\textbf{P4 (25pts)} \cee{if (x{\mbox{<}}={-}20 || x{\mbox{>}}=20) return *(ptrA+1)*4; else return *(ptrB+2)/4;}
\begin{lstlisting}
  cmp   r0,#-20
  ble   then
  cmp   r0,#20
  blt   else_       @ -20<x<20 -> else
then:
  ldr   r0,[r1,#4]  @ *(ptrA+1) = +4 bytes
  lsl   r0,#2       @ * 4
  b     end
else_:
  ldr   r0,[r2,#8]  @ *(ptrB+2) = +8 bytes
  lsr   r0,#2       @ /4 unsigned (use ASR if signed!)
end:
  bx    lr
\end{lstlisting}
\textbf{Trap:} OR short-circuits with first \asm{ble} to \texttt{then}. Ptr offsets in bytes ($k\cdot$sizeof).\end{qp}

% ==========================================================
% [PRIOR SEMESTER PRACTICE] qz-25b material kept for additional practice
\section{Practice Quiz (Prior Sem qz-25b)}

\begin{qp}\textbf{Q5 P1} \cee{if (a{\mbox{>}}=0) return x; else return y;}
\begin{lstlisting}
mp_selector:     @ a,x,y = r0,r1,r2
  cmp   r0,#0
  blt   .else    @ NL on signed
  mov   r0,r1
  bx    lr
.else:
  mov   r0,r2
  bx    lr
\end{lstlisting}\end{qp}

\begin{qp}\textbf{Q5 P2} 5-bit ARM flags, $({-}15){-}({-}16)$: U: $17{-}16{=}1$, no borrow, \textbf{C=1}; S: $1\in[{-}16,15]$, \textbf{V=0}. Binary 0b00001.\\
$16{-}18$: U: needs borrow, \textbf{C=0}; S: $-2\in[-16,15]$, \textbf{V=0}. Binary 0b11110.\end{qp}

\begin{qp}\textbf{Q5 P3} LDR pre/post + scaled. Mem 0x2000\_0000..000F = 0x80..0x87,0x78..0x7F. Two tasks ($[r0,\#4]!$ twice; $[r0],\#8$ then $[r0,r2,lsl \#2]$ with r2=$-1$). Pre-index updates BEFORE; signed reg-offset can be negative.\end{qp}

% ==========================================================
\section{Quiz 6a (flow ctrl)}

\begin{qp}\textbf{Q6a P1} \asm{ldrd/strd} + two \asm{ldrsb [r0,\#4]!}. \asm{ipt[i]=0x70+(i<{<}2)} $\to$ bytes 0x70,0x74,..,0xBC.
After \asm{ldrd r2,r3,[r0,\#8]!}: r2 = little-endian word at ipt+8 = 0x7C78{\textunderscore}7470... wait, bytes are \{0x80,0x84,0x88,0x8C\}, so r2=0x8C{\textunderscore}88{\textunderscore}84{\textunderscore}80.
Then \asm{ldrsb} pre-inc by 4 at 0xB0 $\to$ r2 = 0xFFFFFFB0 (sign-ext MSB=1).\end{qp}

\begin{qp}\textbf{Q6a P2} saturate to $[{-}5,5]$:
\begin{lstlisting}
math_fun_s:          @ x->r0
  cmp     r0,#-5
  blt     .clamp_lo
  cmp     r0,#5
  bgt     .clamp_hi
  bx      lr
.clamp_lo: mov r0,#-5
           bx  lr
.clamp_hi: mov r0,#5
           bx  lr
\end{lstlisting}\end{qp}

\begin{qp}\textbf{Q6a P3} sum of negative elts: \cee{int sum=0; for(i=0;i{\mbox{<}}N;i++) if(arr[i]{\mbox{<}}0) sum+=arr[i];}
\begin{lstlisting}
  mov r2,#0        @ sum
  mov r3,#0        @ i
lp:
  cmp r3,r1
  bge done
  ldr r12,[r0,r3,lsl #2]
  cmp r12,#0
  bge skip
  add r2,r2,r12
skip:
  add r3,#1
  b   lp
done:
  mov r0,r2
  bx  lr
\end{lstlisting}
Standard array-loop template: index, bound-cmp, scaled LDR, predicate, skip, increment, branch.\end{qp}

% ==========================================================
% =========  PART IV: REFERENCE / CHECKLISTS / TRAPS  =====
% ==========================================================

\section{ASM Quick Ref}
\subsection{Immediate loads}
\begin{lstlisting}
ldr   r0,[r1]          @ *r1
ldr   r0,[r1,#4]       @ *(r1+4)
ldr   r0,[r1,#4]!      @ r1+=4; *r1
ldr   r0,[r1],#4       @ *r1; r1+=4
ldrb  r0,[r1]          @ u8 zero-ext
ldrsb r0,[r1]          @ s8 sign-ext
ldrh  r0,[r1]          @ u16 zero-ext
ldrsh r0,[r1]          @ s16 sign-ext
ldrd  r0,r1,[r2,#8]!   @ 64b
\end{lstlisting}

\subsection{Reg-offset / stores}
\begin{lstlisting}
ldr  r0,[r1,r2,lsl #2] @ u32[]
ldrh r0,[r1,r2,lsl #1]
ldrb r0,[r1,r2]
str  r0,[r1,#4]!
strb r0,[r1],#1
strh r0,[r1,#-2]
strd r0,r1,[r2],#8
\end{lstlisting}

\subsection{MOV \& pseudo-LDR (L20)}
\begin{lstlisting}
mov  r0,r1
mov  r0,#42         @ small imm
movw r0,#0xBEEF     @ low 16, zeros high
movt r0,#0xDEAD     @ high 16, preserve low
ldr  r0,=0x12345678 @ pseudo
\end{lstlisting}
\asm{ldr Rd,=expr}: 8-bit $\to$ \asm{mov.w}; 16-bit $\to$ \asm{movw}; 32-bit $\to$ PC-rel \asm{ldr} + \texttt{.word} literal pool.

\subsection{Stack}
\begin{lstlisting}
push {r4-r7,lr}   @ SP-=20; low@low-addr
pop  {r4-r7,pc}   @ restore + return
\end{lstlisting}

% ==========================================================
\section{Directives \& File Header}
\begin{lstlisting}
.section .text
.syntax unified
.thumb
.align 2
.global my_fn
.type   my_fn, %function
my_fn:
  @ ...
  bx  lr
.end
\end{lstlisting}
\texttt{.word}/\texttt{.byte}/\texttt{.asciz} $\to$ literals (TBB table, vector table). \texttt{.weak name} $\to$ placeholder. \texttt{.thumb\_set alias,target} $\to$ symbol alias.

% ==========================================================
\section{Registers \& Exception}
Low R0--R7 (some 16-bit Thumb limited); High R8--R12. R13=SP, R14=LR, R15=PC, R12=IP.

\textbf{xPSR} $=$ \{APSR, IPSR, EPSR\} (overlapping views).

\textbf{Exception auto-stack:} HW pushes \{R0,R1,R2,R3,R12,LR,PC,PSR\} before handler. Handler need NOT re-save these.

% ==========================================================
\section{Instructor Style}
\begin{itemize}
\item Labels: \texttt{<fn>\_lp}, \texttt{\_end}, \texttt{\_then}, \texttt{\_elseif}, \texttt{\_else}.
\item Every ASM fn: (1) C-prototype comment, (2) In/Out/Other reg block, (3) \asm{.global}/\asm{.type}.
\item GNU: lowercase insns; comments \texttt{@}.
\item Return: \asm{bx lr} (leaf) or \asm{pop \{r4,pc\}} (preferred when pushed).
\item Test harness: FUT=Function Under Test, CFN/AFN=C/ASM fn, UTF=Unit Test Fn, NUT=Number Under Test.
\item ITTTT $=$ 4-cond max.
\end{itemize}

% ==========================================================
\section{Glossary}
\begin{tabular}{@{}ll@{}}
\toprule
Term & Meaning \\
\midrule
CCS   & Condition Code Suffix (eq,ne,gt,...) \\
CEX   & Conditional EXec (IT-based) \\
Op2   & ``Flexible 2nd operand'' \\
CP range & Common Positive [0,$2^{N-1}{-}1$] \\
NUT   & Number Under Test (flag analysis) \\
OUT   & Operation Under Test \\
FUT   & Function Under Test \\
Pseudo LDR & \asm{ldr Rd,=expr} \\
Leaf/Stem & no BL / has BL \\
AAPCS & ARM Procedure Call Standard \\
EABI  & Embedded App Binary Interface \\
MAC   & Multiply-Accumulate \\
FD    & Full-descending (stack) \\
\bottomrule
\end{tabular}

% ==========================================================
\section{C$\to$ASM Checklist}
For each C statement:
\begin{enumerate}
\item \textbf{Data type?} $\to$ LDR variant + offset multiplier.
\item \textbf{Inc order?} pre (\texttt{!}) vs post vs none.
\item \textbf{What's modified?} mem value (needs store-back), ptr, or both.
\item \textbf{Signed/unsigned?} $\to$ CCS (GT/LT vs HI/LO), shift (ASR vs LSR), load suf (SB/SH vs B/H).
\item \textbf{Single or branch?} single $\to$ Op2/CEX; multi $\to$ labeled branch.
\item \textbf{Leaf or stem?} any \asm{bl} $\to$ push LR + callee-saved.
\item \textbf{Flags needed?} S-suffix if later branch relies on them.
\end{enumerate}

% ==========================================================
\section{Last-Minute Traps (Exam Gotchas)}
\begin{itemize}
\item \textbf{Sign-ext on byte/half MSbit=1:} LDRSB 0x80$\to$0xFFFFFF80, LDRSH 0xB0A0$\to$0xFFFFB0A0. LDRB/LDRH zero-ext.
\item \textbf{Sign-ext at call site:} \cee{int8\_t 0x81} as arg $\to$ R1=0xFFFFFF81.
\item \textbf{STRD / reg-offset:} neither exists. Don't write \asm{strd rL,rH,[Ra,Ri,LSL\#3]}.
\item \textbf{CBZ/CBNZ in IT:} illegal. Rn must be R0--R7. Forward only, $\le$126B.
\item \textbf{Cond B in IT:} legal ONLY as final slot.
\item \textbf{\#-imm in CMP:} \asm{cmp r0,\#-5} silently becomes \asm{CMN r0,\#5}. Signed CCS still works.
\item \textbf{Prefix \cee{++} on \cee{*p}:} MUST STR/STRB write-back the incremented mem value.
\item \textbf{Ptr arith is scaled:} \cee{p32+1} = +4 bytes, not +1. \cee{(void*)p+1} = +1 byte.
\item \textbf{Flag update requires S-suffix:} ADD $\ne$ ADDS. MUL never sets flags (only MULS sets N,Z).
\item \textbf{UDIV/SDIV never set flags} (no S form). MUL has no C/V.
\item \textbf{ARM sub flag:} C=!Borrow. Opposite of x86.
\item \textbf{Signed vs unsigned divide:} ASR for int/2, LSR for uint/2. ASR on negative number OK; LSR gives garbage for negatives.
\item \textbf{PUSH order:} always ascending reg\# to ascending addr regardless of list syntax.
\item \textbf{LDRD even-odd reg pair rule:} Cortex-M4 allows any pair; some cores require Rt even + Rt+1. Stick with consecutive regs.
\item \textbf{BIC not {\&}{\mbox{\~{}}}:} C has no single ``bit clear'' operator; must spell as \cee{a \& \textasciitilde m}. ASM is \asm{bic Rd,Rn,mask}.
\item \textbf{Op2 shift count:} MUST be immediate (\asm{Rm,LSL \#n}). Standalone LSL/LSR can take Rm for count.
\item \textbf{GNU shifted-Op2:} dest register cannot be omitted --- write \asm{add r0,r1,r1,lsl \#3}, not \asm{add r0,r1,lsl \#3}.
\item \textbf{Q-format mul renorm:} Qm.n $\times$ Qm.n $\to$ Q2m.2n (in 2N-bit); shift right by $n$ to renormalize.
\item \textbf{V=signed, C=unsigned.} Picked via CCS. The HW sets both every arithmetic; you choose which to trust.
\item \textbf{ADR range limited:} \asm{adr r1,table} needs table within PC$\pm$4KB. Use \asm{ldr r1,=table} if far.
\item \textbf{TBB offset scale:} byte table stored as $(case\_label-table)/2$ (Thumb 2-byte alignment).
\end{itemize}

% ==========================================================
% =========  PART V: PATTERN CATALOG (C $\to$ ASM)  =======
% ==========================================================

\section{Pattern Catalog: C $\to$ ASM}

\subsection{Integer arithmetic}
\begin{tabular}{@{}ll@{}}
\toprule
C & ASM \\
\midrule
\cee{c = a + b}        & \asm{add r0,r1,r2} \\
\cee{c = a - b}        & \asm{sub r0,r1,r2} \\
\cee{c = b - a}        & \asm{rsb r0,r1,r2} (Op2$-$Rn) \\
\cee{c = a * b}        & \asm{mul r0,r1,r2} \\
\cee{c = a * b + d}    & \asm{mla r0,r1,r2,r3} \\
\cee{c = d - a * b}    & \asm{mls r0,r1,r2,r3} \\
\cee{c = a / b} unsigned & \asm{udiv r0,r1,r2} \\
\cee{c = a / b} signed   & \asm{sdiv r0,r1,r2} \\
\cee{c = a \% b} unsign. & \asm{udiv rt,r1,r2 / mls r0,rt,r2,r1} \\
\cee{c = -a}           & \asm{rsb r0,r1,\#0} \\
\cee{c = \textasciitilde a}   & \asm{mvn r0,r1} \\
\cee{(int64\_t)a+b}    & \asm{adds rL,rL,\#v / adc rH,rH,\#0} \\
\bottomrule
\end{tabular}

\subsection{Shifts / scaling}
\begin{tabular}{@{}ll@{}}
\toprule
C & ASM \\
\midrule
\cee{c = a {\mbox{<}}{\mbox{<}} n}    & \asm{lsl r0,r1,\#n} \\
\cee{c = (u)a {\mbox{>}}{\mbox{>}} n} & \asm{lsr r0,r1,\#n} \\
\cee{c = (s)a {\mbox{>}}{\mbox{>}} n} & \asm{asr r0,r1,\#n} \\
\cee{c = 2*a}        & \asm{lsl r0,r1,\#1} \\
\cee{c = a/2 signed} & \asm{asr r0,r1,\#1} \\
\cee{c = 3*a}        & \asm{add r0,r1,r1,lsl \#1} \\
\cee{c = 5*a}        & \asm{add r0,r1,r1,lsl \#2} \\
\cee{c = 7*a}        & \asm{rsb r0,r1,r1,lsl \#3} \\
\cee{c = a - a/4 s}  & \asm{sub r0,r1,r1,asr \#2} \\
\cee{c = a + a/8 s}  & \asm{add r0,r1,r1,asr \#3} \\
\bottomrule
\end{tabular}

\subsection{Bit manipulation}
\begin{tabular}{@{}ll@{}}
\toprule
C & ASM \\
\midrule
\cee{c = a \& b}   & \asm{and r0,r1,r2} \\
\cee{c = a | b}    & \asm{orr r0,r1,r2} \\
\cee{c = a \^{} b} & \asm{eor r0,r1,r2} \\
\cee{c = a \& \textasciitilde b}  & \asm{bic r0,r1,r2} \\
\cee{c = a | \textasciitilde b}   & \asm{orn r0,r1,r2} \\
\cee{c |= (1{\mbox{<}}{\mbox{<}}n)}  & \asm{orr r0,r0,\#1,lsl \#n} \\
\cee{c \&= \textasciitilde(1{\mbox{<}}{\mbox{<}}n)} & \asm{bic r0,r0,\#1,lsl \#n} \\
\cee{c \^{}= (1{\mbox{<}}{\mbox{<}}n)}       & \asm{eor r0,r0,\#1,lsl \#n} \\
\cee{bool = c \& (1{\mbox{<}}{\mbox{<}}n)}   & \asm{tst r0,\#1,lsl \#n} \\
\cee{a \&= \textasciitilde(m{\mbox{<}}{\mbox{<}}s); a |= v{\mbox{<}}{\mbox{<}}s} & BIC mask + ORR value \\
\bottomrule
\end{tabular}

\subsection{Memory (pointers, arrays)}
\begin{tabular}{@{}ll@{}}
\toprule
C & ASM \\
\midrule
\cee{v = *p}      & \asm{ldr r0,[r1]} \\
\cee{*p = v}      & \asm{str r0,[r1]} \\
\cee{v = p[k]}    & \asm{ldr r0,[r1,\#4k]} \\
\cee{v = p[i]}    & \asm{ldr r0,[r1,r2,lsl \#2]} \\
\cee{v = *p++}    & \asm{ldr r0,[r1],\#4} \\
\cee{*p++ = v}    & \asm{str r0,[r1],\#4} \\
\cee{v = *++p}    & \asm{ldr r0,[r1,\#4]!} \\
\cee{v = *(int16\_t*)p} & \asm{ldrsh r0,[r1]} \\
\cee{v = *(uint16\_t*)p} & \asm{ldrh r0,[r1]} \\
\cee{v = *(int8\_t*)p}  & \asm{ldrsb r0,[r1]} \\
\cee{v = *(uint8\_t*)p} & \asm{ldrb r0,[r1]} \\
\cee{v = *(int64\_t*)p} & \asm{ldrd r0,r1,[r2]} \\
\cee{p = \&arr[0]} & \asm{ldr r0,=arr} or \asm{adr r0,arr} \\
\bottomrule
\end{tabular}

\subsection{Flow control}
\begin{tabular}{@{}ll@{}}
\toprule
C & ASM \\
\midrule
\cee{if (a==b)}       & cmp / bne end \\
\cee{if (a==0)}       & cbnz r0,skip \\
\cee{while (*p)}      & ldrb + cbz r0,end \\
\cee{if (p==NULL) ret}& cbz r1,end \\
return leaf           & bx lr \\
return stem           & pop \{...,pc\} \\
\cee{f(a,b)}          & mov r0,ra / mov r1,rb / bl f \\
\bottomrule
\end{tabular}

% ==========================================================
\section{Common Recipes}

\subsection{Signed saturate to $[L,H]$}
\begin{lstlisting}
cmp   r0,#L
blt   lo
cmp   r0,#H
bgt   hi
bx    lr
lo:   mov r0,#L / bx lr
hi:   mov r0,#H / bx lr
\end{lstlisting}

\subsection{Array: sum all elements (int)}
\begin{lstlisting}
  mov r2,#0       @ sum
  mov r3,#0       @ i
lp:
  cmp r3,r1       @ vs N
  bge end
  ldr r12,[r0,r3,lsl #2]
  add r2,r2,r12
  add r3,#1
  b   lp
end: mov r0,r2 / bx lr
\end{lstlisting}

\subsection{Array: conditional sum (negatives only)}
\begin{lstlisting}
  mov r2,#0 / mov r3,#0
lp:
  cmp r3,r1
  bge end
  ldr r12,[r0,r3,lsl #2]
  cmp r12,#0
  bge skip
  add r2,r2,r12
skip:
  add r3,#1 / b lp
end: mov r0,r2 / bx lr
\end{lstlisting}

\subsection{Array max (signed)}
\begin{lstlisting}
  ldr r2,[r0]       @ max=arr[0]
  mov r3,#1
lp:
  cmp r3,r1
  bge end
  ldr r12,[r0,r3,lsl #2]
  cmp r12,r2
  ble skip
  mov r2,r12
skip:
  add r3,#1 / b lp
end: mov r0,r2 / bx lr
\end{lstlisting}

\subsection{String length (strlen)}
\begin{lstlisting}
  mov r1,r0
lp:
  ldrb r2,[r1],#1
  cbz  r2,end
  b    lp
end:
  sub  r0,r1,r0
  sub  r0,#1        @ don't count '\0'
  bx   lr
\end{lstlisting}

\subsection{String copy (strcpy)}
\begin{lstlisting}
  @ r0=dst, r1=src
lp:
  ldrb r2,[r1],#1
  strb r2,[r0],#1
  cbnz r2,lp
  bx   lr
\end{lstlisting}

\subsection{Memcpy (byte-wise)}
\begin{lstlisting}
  @ r0=dst, r1=src, r2=n
  mov r3,#0
lp:
  cmp r3,r2
  bge end
  ldrb r12,[r1,r3]
  strb r12,[r0,r3]
  add  r3,#1
  b    lp
end: bx lr
\end{lstlisting}

\subsection{64-bit signed multiply+accumulate MAC}
\begin{lstlisting}
  @ {rH:rL} += rA*rB (signed)
  smlal rL,rH,rA,rB
\end{lstlisting}

\subsection{Q15 multiply (Q15 $\times$ Q15 $\to$ Q15)}
\begin{lstlisting}
  @ r0,r1 are Q15 signed (16-bit in low)
  smull r2,r3,r0,r1   @ Q30 in r3:r2
  lsr   r0,r2,#15     @ renorm Q30 -> Q15
  @ add r0,r0,r3,lsl #17 if high bits matter
\end{lstlisting}

\subsection{Unsigned modulo without UDIV-MLS}
\begin{lstlisting}
  @ fallback: udiv + mul + sub
  udiv r2,r0,r1
  mul  r2,r2,r1
  sub  r0,r0,r2
\end{lstlisting}

\subsection{Pointer difference (elements)}
\begin{lstlisting}
  @ (p - q) of typed pointers
  sub   r0,r0,r1           @ bytes
  asr   r0,r0,#shift       @ /4 for int*, /2 for int16_t*
\end{lstlisting}

\subsection{Assign bit-field (v to bits[s+w-1:s])}
\begin{lstlisting}
  @ r0=target; v in r1; (mask=(1<<w)-1)
  bic r0,r0,#mask,lsl #s
  and r2,r1,#mask           @ trim v
  orr r0,r0,r2,lsl #s
\end{lstlisting}

\subsection{Extract bit-field (bits[s+w-1:s] $\to$ Rd)}
\begin{lstlisting}
  @ UBFX = unsigned bit field extract
  ubfx r0,r1,#s,#w   @ or manual lsr#s+and#mask
  sbfx r0,r1,#s,#w   @ signed (sign-extends)
\end{lstlisting}

% ==========================================================
% =========  PART VI: HW SUMMARY TABLES  ==================
% ==========================================================

\section{HW Coverage Map}

\begin{tabular}{@{}lll@{}}
\toprule
HW & Topic & Key instrs \\
\midrule
HW5 & Q-format encode/decode, Q15 MAC & shifts, mul \\
HW6 & EABI args, 64-bit add/sub, modulo & ADDS/ADC, SUBS/SBC, UDIV \\
HW7 & Shifts, Op2 scaling, bitwise & LSL/LSR/ASR, ORR/BIC/EOR \\
HW8 & NZCV flags, CCS, branch & CMP/CMN, Bcc, SMLAL \\
HW9 & LDR/STR addressing modes & basic/pre/post, reg-offset \\
HW10 & Pointers, LDRD, sign-ext & LDRSB/LDRSH, LDRD \\
HW11 & If/if-else flow control (PL/NL/CEX) & IT blocks, cond B \\
HW12 & Loops, null-term strings & CBZ, post-inc LDRB \\
\bottomrule
\end{tabular}

% ==========================================================
\section{Full HW11 P1 side-by-side}
\cee{if (x{\mbox{>}}y) return (int32\_t)*(++ptrA)*2; else return *(ptrB++)/2;}
R0=x, R1=y, R2=ptrA (int16\_t*), R3=ptrB (int32\_t*).

\subsection{Part 1 (PL)}
\begin{lstlisting}
  cmp r0,r1
  bgt pl_then
pl_else:
  ldr   r0,[r3],#4
  asr   r0,r0,#1
  b     pl_end
pl_then:
  ldrsh r0,[r2,#2]!
  lsl   r0,r0,#1
pl_end:
  bx    lr
\end{lstlisting}

\subsection{Part 2 (NL)}
\begin{lstlisting}
  cmp r0,r1
  ble nl_else
nl_then:
  ldrsh r0,[r2,#2]!
  lsl   r0,r0,#1
  b     nl_end
nl_else:
  ldr   r0,[r3],#4
  asr   r0,r0,#1
nl_end:
  bx    lr
\end{lstlisting}

\subsection{Part 3 (CEX, ITTEE le)}
\begin{lstlisting}
  cmp     r0,r1
  ittee   le
  ldrle   r0,[r3],#4
  asrle   r0,r0,#1
  ldrshgt r0,[r2,#2]!
  lslgt   r0,r0,#1
  bx      lr
\end{lstlisting}

\textbf{Takeaways:} choose CCS suffix=NL-cond in IT header; T-slots match, E-slots inverse. Pre/post-idx LDRs fully legal in IT.

% ==========================================================
\section{Full HW11 P2 piecewise $f(x)$}

\cee{f(x)=\{-10 if x{\mbox{<}}-5; +10 if x{\mbox{>}}=5; 2x else\}}

\subsection{C}
\begin{lstlisting}
int f(int x){
  int y;
  if (x<-5) y=-10;
  else if (x>=5) y=10;
  else y=2*x;
  return y;
}
\end{lstlisting}

\subsection{CMP chain (NL at each arm)}
\begin{lstlisting}
  cmp r0,#-5     @ = CMN r0,#5
  bge elseif
  ldr r0,=-10
  b   end
elseif:
  cmp r0,#5
  blt else_blk
  ldr r0,=10
  b   end
else_blk:
  lsl r0,r0,#1   @ 2*x
end:
  bx  lr
\end{lstlisting}

\subsection{CEX (chained IT + cond B-as-last-slot)}
\begin{lstlisting}
  cmp   r0,#-5
  itt   lt
  ldrlt r1,=-10  @ T1
  blt   end      @ T2 (last slot!)
  cmp   r0,#5
  ite   ge
  ldrge r1,=10
  lsllt r1,r0,#1 @ uses orig x in r0
end:
  mov   r0,r1
  bx    lr
\end{lstlisting}

\textbf{Takeaway:} stage result in R1 so R0 (holding x) survives until 2nd CMP. Cond B in IT is legal \textbf{only as last slot}.

% ==========================================================
\section{Full HW11 P3 compound OR}
\cee{if (x{\mbox{<}}=20 || x{\mbox{>}}=25) ... else ...}
R0=x, R1=ptrA (uint16\_t*), R2=ptrB (uint32\_t*). Unsigned $/4\Rightarrow$ LSR.

\subsection{CMP short-circuit}
\begin{lstlisting}
  cmp r0,#20
  ble then
  cmp r0,#25
  bge then
  b   else_
then:
  ldrh r0,[r1,#2]!
  lsl  r0,r0,#2
  b    end
else_:
  ldr  r0,[r2,#4]!
  lsr  r0,r0,#2
end:
  bx   lr
\end{lstlisting}

\subsection{CEX (1st test plain B, 2nd test ITTEE)}
\begin{lstlisting}
  cmp r0,#20
  bgt check25
  ldrh r0,[r1,#2]!
  lsl  r0,r0,#2
  bx   lr
check25:
  cmp r0,#25
  ittee  ge
  ldrhge r0,[r1,#2]!
  lslge  r0,r0,#2
  ldrlt  r0,[r2,#4]!
  lsrlt  r0,r0,#2
  bx     lr
\end{lstlisting}

\textbf{Takeaway:} 1st of two OR sub-tests can't pack into IT (cond B only at end); handle with plain \asm{bgt}, then ITTEE the tail.

% ==========================================================
\section{Full HW12 P2 string copy (upper-case)}

\subsection{C (while with priming read)}
\begin{lstlisting}
void mp_str_cpy_2_caps_while_c(char *d, char *s){
  char ch = *s++;
  while (ch){
    ch -= 0x20;
    *d++ = ch;
    ch = *s++;     @ duplicate read
  }
  *d++ = ch;       @ writes '\0'
}
\end{lstlisting}

\subsection{Part 1 ASM}
\begin{lstlisting}
mp_str_cpy_2_caps_while_s:
  ldrb r2,[r1],#1  @ priming
lp:
  cbz  r2,end
  sub  r2,r2,#0x20
  strb r2,[r0],#1
  ldrb r2,[r1],#1  @ in-loop DUP
  b    lp
end:
  strb r2,[r0],#1  @ '\0' (r2=0 via CBZ)
  bx   lr
\end{lstlisting}

\subsection{Part 2 C (while(1)+break)}
\begin{lstlisting}
void mp_str_cpy_2_caps_while_brk_c(char *d, char *s){
  while (1){
    char ch = *s++;
    if (ch == 0) { *d++ = ch; break; }
    ch -= 0x20;
    *d++ = ch;
  }
}
\end{lstlisting}

\subsection{Part 3 ASM (single LDRB text)}
\begin{lstlisting}
mp_str_cpy_2_caps_while_brk_s:
lp:
  ldrb r2,[r1],#1
  cbz  r2,end
  sub  r2,r2,#0x20
  strb r2,[r0],#1
  b    lp
end:
  strb r2,[r0],#1
  bx   lr
\end{lstlisting}

\textbf{Takeaway:} priming-read style $\to$ two LDRB texts (easy error source). while(1)+break $\to$ one LDRB. Same number of dynamic LDRBs (N+1).

% ==========================================================
\section{Test Framework / Naming}
\begin{itemize}
\item \textbf{FUT} = Function Under Test.
\item \textbf{CFN} = C Function impl. \textbf{AFN} = ASM Function impl.
\item \textbf{UTF} = Unit Test Function.
\item \textbf{NUT} = Number Under Test (flag-bit analysis).
\item \textbf{OUT} = Operation Under Test.
\item \textbf{CP} = Common Positive range $[0, 2^{N-1}{-}1]$ where signed/unsigned agree.
\item Typical name: \texttt{mp\_<op>\_<type>\_<c|s>} e.g.\ \texttt{mp\_add\_64bit\_s}.
\item Each fn has: (1) \texttt{/* C-prototype comment */}, (2) In/Out/Other reg block, (3) \asm{.global} + \asm{.type \%function}.
\end{itemize}

% ==========================================================
\section{GPIO Quick Snippets (L6--7)}

These appear on earlier labs; unlikely on the final but useful if asked:

\subsection{Direct ODR access}
\begin{lstlisting}
@ set PA5 high via ODR OR (read-modify-write)
ldr r0,=GPIOA_ODR
ldr r1,[r0]
orr r1,r1,#(1<<5)
str r1,[r0]
\end{lstlisting}

\subsection{Atomic BSRR}
\begin{lstlisting}
@ set PA5: write 1<<5 to BSRR[15:0]
ldr r0,=GPIOA_BSRR
mov r1,#(1<<5)
str r1,[r0]

@ reset PA5: write 1<<(5+16) to BSRR[31:16]
mov r1,#(1<<(5+16))
str r1,[r0]
\end{lstlisting}

\textbf{BSRR advantage:} single atomic write, no read-modify-write race with ISR modifying other bits of same port.

% ==========================================================
\section{Exception Handler Entry (L8)}
HW auto-stacks \{R0, R1, R2, R3, R12, LR, PC, PSR\} on current stack. Handler inherits LR=\texttt{EXC\_RETURN} magic value (\texttt{0xFFFFFFFx}) --- BX LR returns from the exception.

\begin{itemize}
\item Handler need NOT re-save R0--R3, R12, LR, PC, PSR.
\item Handler DOES need to save R4--R11 if it uses them.
\item EXC\_RETURN picks between MSP/PSP and thread/handler mode.
\end{itemize}

% ==========================================================
\section{Memory Regions (F412 / G431)}
\begin{tabular}{@{}lll@{}}
\toprule
Region & Start & Use \\
\midrule
Flash       & 0x08000000 & code, \texttt{.rodata}, literal pool \\
SRAM        & 0x20000000 & \texttt{.data}, \texttt{.bss}, stack \\
Peripheral  & 0x40000000 & GPIO/UART/TIM regs \\
System      & 0xE0000000 & NVIC, SysTick \\
\bottomrule
\end{tabular}
Vector table at 0x08000000 by default: SP init (word 0), reset handler (word 1), then NMI, HardFault, ... IRQs.

% ==========================================================
\section{Endianness Drills}
ARM Cortex-M4 = \textbf{little-endian}. Byte at lowest addr = LSB of the multi-byte word.

\begin{tabular}{@{}ll@{}}
\toprule
Bytes @addr & 32-bit word \\
\midrule
0x01,0x02,0x03,0x04 & 0x04030201 \\
0x80,0x00,0x00,0x00 & 0x00000080 \\
0xFF,0xFE,0xFD,0xFC & 0xFCFDFEFF \\
\bottomrule
\end{tabular}

\textbf{LDR alignment:} word LDR needs 4-byte aligned addr; LDRH needs 2-byte aligned; LDRB any. Unaligned LDR may fault (Cortex-M4 allows, but slow).

% ==========================================================
\section{Numeric Ranges}
\begin{tabular}{@{}lll@{}}
\toprule
Type & Min & Max \\
\midrule
uint8\_t  & 0 & 255 \\
int8\_t   & $-$128 & 127 \\
uint16\_t & 0 & 65535 \\
int16\_t  & $-$32768 & 32767 \\
uint32\_t & 0 & $2^{32}{-}1$ \\
int32\_t  & $-2^{31}$ & $2^{31}{-}1$ \\
int64\_t  & $-2^{63}$ & $2^{63}{-}1$ \\
\bottomrule
\end{tabular}

\textbf{TC encode:} neg $\to$ add $2^N$ (wrap). \textbf{TC decode:} if MSB=1, subtract $2^N$.

\textbf{Hex $\leftrightarrow$ dec ($8$ bit signed):}
\begin{tabular}{@{}lll@{}}
\toprule
Hex & TC dec & sign bit \\
\midrule
0x00 & 0      & + \\
0x7F & 127    & + (max+) \\
0x80 & -128   & - (min-) \\
0xFF & -1     & - \\
\bottomrule
\end{tabular}

% ==========================================================
\section{Full ASM for HW10 P1 (4 lines)}

\cee{uint8\_t *p8a = 0x20000100;  uint8\_t *p8b = 0x2000010C;}

Lines:
\begin{lstlisting}
@ Line 4: *p8b++ = *(p8a + 2);
  ldrb r2,[r0,#2]
  strb r2,[r1],#1

@ Line 5: *p8b++ = ++*(p8a + 4);
  ldrb r2,[r0,#4]
  add  r2,r2,#1
  strb r2,[r0,#4]     @ write-back!
  strb r2,[r1],#1

@ Line 6: *p8b++ = ++*p8a++;
  ldrb r2,[r0],#1
  add  r2,r2,#1
  strb r2,[r0,#-1]    @ back to old addr
  strb r2,[r1],#1

@ Line 7: *p8b++ = *p8a++;
  ldrb r2,[r0],#1
  strb r2,[r1],#1
\end{lstlisting}

Prefix \cee{++} on a \cee{*}-deref $\Rightarrow$ STR write-back after ADD. Postfix on pointer $\Rightarrow$ post-index.

% ==========================================================
\section{Full HW9 P1 ASM}

16 bytes 0x00..0x0F @ 0x20000000. Little-endian word layout:

\begin{tabular}{@{}ll@{}}
\toprule
Addr & Word \\
\midrule
0x20000000 & 0x03020100 \\
0x20000004 & 0x07060504 \\
0x20000008 & 0x0B0A0908 \\
0x2000000C & 0x0F0E0D0C \\
\bottomrule
\end{tabular}

\begin{lstlisting}
ldr r4,=0x20000000    @ r4 = 0x20000000
ldr r0,[r4,#4]        @ r0 = 0x07060504
                       @ r4 unchanged
ldr r1,[r4,#2]!       @ r4 = 0x20000002
                       @ r1 = 0x05040302
ldr r2,[r4],#4        @ r2 = 0x05040302
                       @ r4 = 0x20000006
ldr r3,[r4]           @ r3 = 0x09080706
\end{lstlisting}

Critical: \texttt{!}-pre-idx updates \texttt{Ra} BEFORE load; post-idx AFTER.

% ==========================================================
\section{Common Exam Traps --- Extended}
\begin{itemize}
\item \textbf{Misreading LDR alignment:} \asm{ldr r0,[r4,\#2]!} loads FOUR bytes starting at $r_4+2$ (so bytes 2..5), NOT two bytes.
\item \textbf{LDRD register pair is consecutive:} \asm{ldrd r2,r3,...} --- r3 must be r2+1. \asm{ldrd r2,r5,...} is illegal.
\item \textbf{STRB silently discards upper 24 bits.} Writing \cee{0x12345678} via STRB stores only \cee{0x78}.
\item \textbf{ASR vs LSR on boundary values:} \asm{asr r0,r1,\#1} where r1=0x80000000 gives 0xC0000000 (sign fill); \asm{lsr} gives 0x40000000.
\item \textbf{CMP immediate range:} 12-bit modified immediate. Big constants must be loaded via \asm{ldr r12,=big}, then CMP reg-to-reg.
\item \textbf{MOVS vs MOV:} MOVS sets flags. \asm{movs r0,\#0} clears Z=1, N=0. Handy for zeroing AND testing.
\item \textbf{UDIV/SDIV by zero:} returns 0 on Cortex-M4 (no exception by default). Check divisor first if needed.
\item \textbf{Integer overflow on MUL:} 32-bit MUL truncates; bits above 31 silently lost. Use SMULL/UMULL to preserve.
\item \textbf{EABI stack alignment:} before a BL, SP must be 8-byte aligned. Push odd \# of regs + align dummy, or push pairs.
\item \textbf{CBZ target must be AHEAD in memory.} Backward CBZ illegal.
\item \textbf{IT block terminates early} if interrupted by LDR literal pool or label.
\item \textbf{Cortex-M4 does NOT support ARM mode} --- Thumb-2 only. Forgetting \texttt{.thumb} directive breaks everything.
\item \textbf{Tail-call optimization:} \asm{b func} (not bl) as the last instruction of a leaf-after-BL chain re-uses caller's LR --- saves a pop/push pair.
\item \textbf{Self-modifying stack via PUSH+POP order:} PUSH preserves list regardless of order; POP restores in the same layout. Don't rely on list-syntax ordering.
\end{itemize}

% ==========================================================
\section{Worked Flag-Compute Drill}

\textbf{Example:} 8-bit sys ($C_N=256$), compute flags for \asm{r0 + r1} where r0=0xA3 (163), r1=0x7F (127).

\textbf{Unsigned:} $163+127 = 290$. $290>U_{\max}=255$ $\Rightarrow$ $290-256=34$. \textbf{C=1}.

\textbf{Signed:} $r_0$ = 0xA3 has MSB=1 so value $=163-256=-93$; $r_1$ = 0x7F = +127. $-93+127 = 34 \in [-128,127]$ $\Rightarrow$ \textbf{V=0}.

\textbf{Result binary:} $34 = $ 0b00100010. N=0 (MSB). Z=0.

\textbf{Final flags:} N=0, Z=0, C=1, V=0.

---

\textbf{Example 2:} 8-bit $r_0=0x7F$ (+127) $+$ $r_1=0x01$. Unsigned: $128<256$, no carry, \textbf{C=0}. Signed: $127+1=128>127$ $\Rightarrow$ \textbf{V=1} (corrected to $-128$). Binary 0b10000000. N=1, Z=0.

\textbf{Subtraction drill:} $r_0-r_1$ where $r_0=0x20$ (32), $r_1=0x30$ (48). Unsigned $32-48$ needs borrow $\Rightarrow$ \textbf{C=0}. $32-48=-16 \in [-128,127]$ $\Rightarrow$ \textbf{V=0}. Result $-16 = $ 0xF0. N=1, Z=0.

% ==========================================================
\section{Thumb-2 Reserved / Special Names}
\begin{tabular}{@{}ll@{}}
\toprule
Name & Alias \\
\midrule
R13 & SP  (stack ptr) \\
R14 & LR  (link reg) \\
R15 & PC  (program counter) \\
R12 & IP  (intra-proc scratch) \\
\bottomrule
\end{tabular}

\textbf{Don't use as general regs:} SP (push/pop only), PC (branching only). LR is free-ish but must preserve if stem.

\textbf{``Low'' registers (R0--R7):} available to most 16-bit Thumb insns. Some insns (e.g.\ CBZ) restrict to these.

% ==========================================================
\section{Two's Complement Speed Sheet}
\textbf{Negate} (2's comp): invert all bits, add 1.

\begin{tabular}{@{}lll@{}}
\toprule
8-bit hex & dec & binary \\
\midrule
0x01 & +1   & 0000 0001 \\
0x7F & +127 & 0111 1111 \\
0x80 & -128 & 1000 0000 \\
0xFF & -1   & 1111 1111 \\
0xC0 & -64  & 1100 0000 \\
0xA3 & -93  & 1010 0011 \\
\bottomrule
\end{tabular}

\textbf{Quick rule:} 8-bit TC dec $=$ (unsigned dec) $- 256$ when MSB=1.

\textbf{Sign-ext byte to 32-bit:} if top bit of 8-bit value is 1, set high 24 bits to 1s: \cee{0x81 $\to$ 0xFFFFFF81}. Otherwise 0 fill: \cee{0x7E $\to$ 0x0000007E}.

% ==========================================================
\section{Stack Drawings}

\textbf{Initial state} --- SP = 0x20001020:
\begin{verbatim}
0x20001020  (top, nothing pushed yet)
\end{verbatim}

\textbf{After} \asm{push \{r4,r8\}} (2 regs, 8B):
\begin{verbatim}
0x20001018  R4   <-- SP
0x2000101C  R8
0x20001020  (was top)
\end{verbatim}

Regs stored at LOW addr to HIGH addr in \textbf{ascending reg order}: R4 low, R8 high.

\textbf{After} \asm{pop \{r4,r8\}}: SP restored to 0x20001020; slots unchanged in memory (not erased, just forgotten).

% ==========================================================
\section{Condition-Flag Quick Lookup}

After \asm{CMP r0, r1} (i.e.\ \asm{r0 - r1}):

\begin{tabular}{@{}lll@{}}
\toprule
If \ldots & then \\
\midrule
r0 == r1 & Z=1, C=1, N=0, V=0 \\
r0 $>$ r1 (u) & Z=0, C=1 \\
r0 $<$ r1 (u) & C=0 \\
r0 $>$ r1 (s) & Z=0, V==N \\
r0 $<$ r1 (s) & V $\ne$ N \\
r0 $=$ 0, r1 $\ne$ 0 (s) & N depends on r1 sign \\
\bottomrule
\end{tabular}

\textbf{Branch chooser:} if C gives \cee{==}/\cee{!=}, just use EQ/NE --- works for both signed and unsigned.

% [TRIM-Part1] Calling-Convention Templates removed (covered by Function Calls section)
% [MOVE-Part2] Operator Precedence moved to Part 2 (C reference section)

% ==========================================================
\section{HW5 Walkthrough (all 5 problems)}

\subsection{P1 (8pts) Encode 2.75 in UQ3.5}
$N=8, n=5, \mathrm{scale}=2^5=32$. $I=2.75 \cdot 32 = 88 = $\textbf{0x58}.

\subsection{P2 (16pts) Decode Q3.4 codes}
\textbf{(a)} 0x77 = 119, MSB=0 (positive), $f = 119/16 = $ \textbf{+7.4375}.\\
\textbf{(b)} 0xAA = 170, MSB=1 (negative), $I = 170-256=-86$, $f = -86/16 = $ \textbf{-5.375}.

\subsection{P3 (10pts) Decode Q2.5 code 0b10011001}
$U=153$, MSB=1, $I = 153-256=-103$, $f = -103/32 =$ \textbf{-3.21875}.

\subsection{P4 (32pts) Encode in Q3.4}
\textbf{(a)} $0.5 \cdot 16 = 8 \to$ \textbf{0x08}\\
\textbf{(b)} $6.73 \cdot 16 = 107.68 \to 108 \to$ \textbf{0x6C} (err 0.02)\\
\textbf{(c)} $-2.25 \cdot 16 = -36, +256 = 220 \to$ \textbf{0xDC}\\
\textbf{(d)} $-4.5 \cdot 16 = -72, +256 = 184 \to$ \textbf{0xB8}

\subsection{P5 (44pts) Q15 MAC: $f_1 f_2 + f_3 f_4$}
Encode (multiply by $2^{15}$):
\begin{tabular}{@{}llll@{}}
$f$ & $\times 2^{15}$ & hex & $0.5$ \\
0.5    & 16384 & 0x4000 \\
0.25   &  8192 & 0x2000 \\
0.125  &  4096 & 0x1000 \\
0.0625 &  2048 & 0x0800 \\
\end{tabular}

Multiply Q15 (Q15 $\times$ Q15 $=$ Q30; renorm by $\gg 15$):\\
$0.5 \cdot 0.25$: $0x4000 \cdot 0x2000 = 0x08000000 \gg 15 = $ \textbf{0x1000} (= 0.125)\\
$0.125 \cdot 0.0625$: $0x1000 \cdot 0x0800 = 0x00800000 \gg 15 = $ \textbf{0x0100} (= 0.0078125)

Add: $0x1000 + 0x0100 =$ \textbf{0x1100} $= 4352/32768 \approx$ \textbf{0.1328125}.

% ==========================================================
\section{HW6 Walkthrough (7 problems)}

\subsection{P1 PUSH on full-descending stack}
SP=0x20001020; \asm{push \{r7,r5,r8,r6\}}; \cee{Ri=(1<{<}i)+(1<{<}2i)}.\\
P1.1 SP after = 0x20001020$-$16 = \textbf{0x20001010}.\\
P1.2 R5 = $32+1024=$ \textbf{0x00000420}.\\
P1.3 Storage low$\to$high: R5@..1010, R6@..1014, R7@..1018, R8@..101C. Word@..1014 = R6 = $64+4096=$ \textbf{0x00001040}.\\
P1.4 Byte@0x20001014 = LSB of R6 = \textbf{0x40} (LE).

\subsection{P2 EABI unsigned promote}
Call \cee{fn1((uint8\_t)0xF, (uint8\_t)0xFF, 0xFFF, 0xFFFF)}; fn1 takes \cee{uint32\_t}.\\
R0=0x0000000F, R1=0x000000FF, R2=0x00000FFF, R3=0x0000FFFF (zero-ext).

\subsection{P3 EABI signed promote}
Call \cee{fn2((int8\_t)0x0F, (int8\_t)0xFF)}; fn2 takes \cee{int32\_t}.\\
R0 = 0x0000000F (MSB=0).\\
R1: 0xFF = $-1$ as int8; sign-ext $\to$ \textbf{0xFFFFFFFF}.

\subsection{P4 Q15.16 return}
Encode 3.25: $I = 3.25 \cdot 2^{16} = 3 \cdot 0x10000 + 0x4000 = $ \textbf{0x00034000}.\\
32-bit return $\to$ R0.

\subsection{P5 64-bit ADD}
Convention: R0=lo(A), R1=hi(A), R2=lo(B), R3=hi(B).
\begin{lstlisting}
mp_add_64bit_s:
  adds r0,r2     @ lo+lo, sets C
  adc  r1,r3     @ hi+hi+C
  bx   lr
\end{lstlisting}
T1: 0x88000000+0x89000000 = 0x111000000 $\Rightarrow$ lo=0x11000000, C=1; hi=3+2+1=6.\\
T2: 0x78000000+0x79000000 = 0xF1000000 $\Rightarrow$ lo=0xF1000000, C=0; hi=3+2+0=5.

\subsection{P6 64-bit SUB}
\begin{lstlisting}
mp_sub_64bit_s:
  subs r0,r2     @ lo-lo, C=!borrow
  sbc  r1,r3     @ hi-hi-(1-C)
  bx   lr
\end{lstlisting}
T1: 0x70000000$-$0x80000000 needs borrow $\Rightarrow$ lo=0xF0000000, C=0; hi=0x38$-$0x02$-$1=0x35.\\
T2: 0x80000000$-$0x80000000=0, no borrow $\Rightarrow$ C=1; hi=0x38$-$0x28$-$0=0x10.

\subsection{P7 UDIV+MUL modulo}
$A\%B = A - (A/B) \cdot B$.
\begin{lstlisting}
mp_umod_32bit_s:
  udiv r2,r0,r1
  mul  r2,r2,r1
  sub  r0,r0,r2
  bx   lr
\end{lstlisting}
A=17, B=5: $17/5=3$, $3\cdot5=15$, $17-15=$\textbf{2}.

\textbf{Alt (1 fewer instr):} \asm{udiv r2,r0,r1 / mls r0,r2,r1,r0} (\asm{mls Rd,Rn,Rm,Ra} $\to$ \cee{Rd=Ra-Rn*Rm}).

% ==========================================================
\section{HW7 Walkthrough (6 problems)}

\subsection{P1 Q15.16 from SMULL}
After \asm{smull r1,r0,r0,r1} (R0:HW, R1:LW), extract bits b47..b16 to put result in R0:
\begin{lstlisting}
smull r1,r0,r0,r1
lsl   r0,#16          @ HW<<(target_MSB-31)
add   r0,r0,r1,lsr #16
\end{lstlisting}
Pattern: \asm{lsl} HW by ($\text{target MSB}-31$); \asm{lsr} LW by $n$. Q31: lsl1+lsr31. Q15.16: lsl16+lsr16. Q16.15: lsl17+lsr15.

\subsection{P2 LSLS/LSRS and C}
R0=0xDD = 0b11011101.\\
\asm{lsls r1,r0,\#1} $\to$ R1=0x1BA, C $=$ b7 of R0 $=$ 1.\\
\asm{lsrs r2,r0,\#3} $\to$ R2=0x1B, C $=$ b2 of R0 $=$ 1.

\subsection{P3 One-instr scaling (signed)}
\begin{tabular}{@{}ll@{}}
\toprule
$A$ & ASM \\
\midrule
$16B$    & \asm{lsl r0,r1,\#4} \\
$17B$    & \asm{add r0,r1,r1,lsl \#4} \\
$15B$    & \asm{rsb r0,r1,r1,lsl \#4} \\
$B/16$   & \asm{asr r0,r1,\#4} \\
$17B/16$ & \asm{add r0,r1,r1,asr \#4} \\
$15B/16$ & \asm{sub r0,r1,r1,asr \#4} \\
\bottomrule
\end{tabular}

\subsection{P4 BIC+ORR replace nibble}
$0xH_3 H_2 H_1 H_0 \to 0xH_3 H_2 A H_0$ (set bits[7:4]=0xA):
\begin{lstlisting}
ldr r1,=15
bic r0,r0,r1,lsl #4
ldr r2,=10
orr r0,r0,r2,lsl #4
\end{lstlisting}

\subsection{P5 Hand-eval BIC/ORR/EOR}
R0=0xAD, R1=0x03:\\
\asm{bic r3,r0,r1} = 0xAD\,\&\,$\sim$0x03 = 0xAD\,\&\,0xFC = \textbf{0xAC}\\
\asm{orr r4,r3,r1} = 0xAC | 0x03 = \textbf{0xAF}\\
\asm{eor r5,r0,r1} = 0xAD $\oplus$ 0x03 = \textbf{0xAE}

\subsection{P6 Insert digit at end}
\textbf{(b)} Shift left + insert 3 at bottom:
\begin{lstlisting}
lsl r0,#4
orr r0,#3
\end{lstlisting}
\textbf{(c)} Shift right + insert 7 at top (16-bit val):
\begin{lstlisting}
lsr r0,#4
ldr r1,=7
lsl r1,#12
orr r0,r1
\end{lstlisting}

% ==========================================================
\section{Worked CCS Drill}

\textbf{Setup:} 6-bit system $C_N=64$. After \asm{cmp r0,r1} where $r_0=22, r_1=20$ (or equivalently $r_0=-20, r_1=-22$):

Compute $r_0 - r_1 = 2$ in 6-bit:
\begin{itemize}
\item \textbf{Unsigned:} no borrow $\Rightarrow$ \textbf{C=1}.
\item \textbf{Signed:} $2 \in [-32, 31]$, no overflow $\Rightarrow$ \textbf{V=0}.
\item N=0, Z=0.
\end{itemize}

\textbf{Branch evaluation} (all CCS):

\begin{tabular}{@{}llll@{}}
\toprule
CCS & Logic & Value & Branch? \\
\midrule
EQ & Z=1 & 0 & no \\
NE & Z=0 & 1 & yes \\
GT & V=N and Z=0 & 1 & yes \\
GE & V=N & 1 & yes \\
LT & V$\ne$N & 0 & no \\
LE & V$\ne$N or Z=1 & 0 & no \\
HI & C=1 and Z=0 & 1 & yes \\
HS & C=1 & 1 & yes \\
LO & C=0 & 0 & no \\
LS & C=0 or Z=1 & 0 & no \\
\bottomrule
\end{tabular}

\textbf{Same flags from a different CMP:} $-20-(-22)$: in unsigned-mapped form $44-42=2$. Same flags. CCS picks same way.

% ==========================================================
\section{Full HW9 P3 Tasks}

\subsection{Task 1 (basic offsets)}
\begin{lstlisting}
ldr r2,[r0,#4]     @ r2 = ipt+4 word = 0x70605040
ldr r3,[r1,#0]     @ store via opt
@ printout: 0x70605040
ldr r2,[r0,#0xC]
@ printout: 0xF0E0D0C0
\end{lstlisting}
After: R0 unchanged, R1 unchanged.

\subsection{Task 2 (pre-index $\#4!$)}
\begin{lstlisting}
ldr r2,[r0,#4]!    @ r0+=4 first
                    @ r0=0x...14, r2=0x70605040
ldr r2,[r0,#4]!    @ r0=0x...18, r2=0xB0A09080
\end{lstlisting}
After: R0=0x20001018.

\subsection{Task 3 (post-index $\#4$)}
\begin{lstlisting}
ldr r2,[r0],#4     @ r2=0x30201000, then r0+=4
ldr r2,[r0],#4     @ r2=0x70605040, then r0+=4
\end{lstlisting}
After: R0=0x20001018.

\subsection{Task 4 (register-offset, r2=1)}
\begin{lstlisting}
ldr r3,[r0,r2,lsl #2]   @ r2=1: addr=ipt+4
                          @ r3=0x70605040
add r2,#1
ldr r3,[r0,r2,lsl #2]   @ r2=2: addr=ipt+8
                          @ r3=0xB0A09080
add r2,#1               @ r2=3
\end{lstlisting}
After: R0 unchanged, R2=3, R3=0xB0A09080.

% ==========================================================
% [TRIM-Part1] Bit-Twiddle Idioms moved to Part 2 reference section

% ==========================================================
\section{Common ASM-File Skeleton}

\begin{lstlisting}
@ ============================================
@ File: mp_my_func.s
@ Purpose: <brief>
@ Author: <name>
@ Prototype: int mp_my_func(int x, int y);
@ ============================================
.section .text
.syntax unified
.thumb
.align 2

@ -----------------------------
@ In:  R0 = x  (int32_t)
@      R1 = y  (int32_t)
@ Out: R0 = result
@ Other: R12 scratch
@ -----------------------------
.global mp_my_func_s
.type   mp_my_func_s, %function
mp_my_func_s:
  push {lr}
  @ ... body ...
  pop  {pc}
.end
\end{lstlisting}

% ==========================================================
\section{Number-Format Speed Tables}

\subsection{8-bit signed (TC) reference}
\begin{tabular}{@{}lllll@{}}
\toprule
Hex & U & S & Hex & U / S \\
\midrule
0x00 & 0 & 0       & 0x80 & 128 / -128 \\
0x01 & 1 & +1      & 0x81 & 129 / -127 \\
0x10 & 16 & +16    & 0xC0 & 192 / -64 \\
0x40 & 64 & +64    & 0xE0 & 224 / -32 \\
0x7F & 127 & +127  & 0xFF & 255 / -1 \\
\bottomrule
\end{tabular}

\subsection{16-bit signed boundaries}
\begin{tabular}{@{}lll@{}}
\toprule
Hex & U & S \\
\midrule
0x0000 & 0     & 0 \\
0x7FFF & 32767 & +32767 \\
0x8000 & 32768 & -32768 \\
0xFFFF & 65535 & -1 \\
\bottomrule
\end{tabular}

\subsection{Q15 example values}
$0.5 = $ 0x4000, $0.25 = $ 0x2000, $0.125 = $ 0x1000, $0.0625 = $ 0x0800, $-0.5 = $ 0xC000, $-1.0 = $ 0x8000 (saturates), $+0.99997 = $ 0x7FFF (max).

% ==========================================================
\section{Quick Bcc Decoder}
\begin{tabular}{@{}lll@{}}
\toprule
Bcc & ``Branch if'' & Flags \\
\midrule
B    & always   & --- \\
BEQ  & ==       & Z=1 \\
BNE  & !=       & Z=0 \\
BHS/BCS & u $\ge$ & C=1 \\
BLO/BCC & u $<$  & C=0 \\
BMI  & N=1 (negative result) & N=1 \\
BPL  & N=0 (zero or pos)     & N=0 \\
BVS  & V=1 (signed overflow) & V=1 \\
BVC  & V=0  & V=0 \\
BHI  & u $>$  & C=1 and Z=0 \\
BLS  & u $\le$ & C=0 or Z=1 \\
BGE  & s $\ge$ & V=N \\
BLT  & s $<$   & V $\ne$ N \\
BGT  & s $>$   & V=N and Z=0 \\
BLE  & s $\le$ & V$\ne$N or Z=1 \\
\bottomrule
\end{tabular}

% ==========================================================
\section{HW8 Walkthrough (NZCV+CCS)}

\subsection{P1 (42pts) Compute NZCV by hand (5-bit)}
$C_N=32, U_{\max}=31, S_{\min}=-16, S_{\max}=15$.

\textbf{Procedure} per add/sub:
\begin{enumerate}
\item Convert negatives to unsigned via $+C_N$.
\item Compute as if unsigned in $[0, C_N)$. If result wraps $\Rightarrow$ \textbf{C} (add: C=1; sub: borrow C=0).
\item Treat the SAME binary as signed. If result $\notin [S_{\min}, S_{\max}]$, \textbf{V}=1, correct by $\pm C_N$.
\item N=MSB of binary; Z=(value$=$0).
\end{enumerate}

\textbf{Drill 1:} $({-}15)+({-}5)$.\\
U: $17 + 27 = 44 \to 44-32=12$, \textbf{C=1}.\\
S: $-15+(-5)=-20 < -16$, $+32=12$, \textbf{V=1}.\\
Binary 0b01100, N=0, Z=0.

\textbf{Drill 2:} $11-18$.\\
U: $11-18=-7 \to +32=25$, borrow $\Rightarrow$ \textbf{C=0}.\\
S: $11-18=-7 \in [-16,15]$, \textbf{V=0}.\\
Binary 0b11001, N=1, Z=0.

\textbf{Drill 3:} $7+9$ (5-bit).\\
U: $7+9=16 \in [0,31]$, no carry, \textbf{C=0}.\\
S: $7+9=16 > 15$, \textbf{V=1}, $-32=-16$.\\
Binary 0b10000, N=1, Z=0. (Signed reads as $-16$, unsigned reads as $16$.)

\subsection{P2 CCS table (after CMP)}
Compute flags from \asm{cmp r0,r1} as $r_0 - r_1$, then evaluate every CCS. Equal-magnitude pairs (e.g.\ $22-20$ vs $-20-(-22)$) give identical flags because $(-20)+C_N - ((-22)+C_N) = 22-20 = 2$.

\textbf{Trick:} CCS is \textbf{flag-driven}, not value-driven. Two CMPs that produce identical flags branch identically.

\subsection{P3 max() via NL branch}
\cee{int max(int a,int b)\{ return a{\mbox{>}}=b ? a : b; \}}
\begin{lstlisting}
max:
  cmp r0,r1
  blt .else      @ NL: branch on !(a>=b)
  bx  lr          @ then: r0 already = a
.else:
  mov r0,r1       @ r0 = b
  bx  lr
\end{lstlisting}

\subsection{P4 64-bit sum-of-squares with SMLAL}
\cee{int64\_t sum=0; for(i=s; i{\mbox{<}}=e; i++) sum += i*i;}
\begin{lstlisting}
  ldr r2,=0       @ sum_lo
  ldr r3,=0       @ sum_hi
.lp:
  cmp r0,r1
  bgt .end
  smlal r2,r3,r0,r0  @ {hi:lo}+=i*i
  add   r0,#1
  b     .lp
.end:
  mov   r0,r2     @ lo -> R0
  mov   r1,r3     @ hi -> R1
  bx    lr
\end{lstlisting}

\textbf{Key:} \asm{SMLAL Rlo,Rhi,Rn,Rm}: signed 32$\times$32$\to$64-bit accumulate. Pair with EABI 64-bit return (R0=lo, R1=hi).

% ==========================================================
\section{Worked Inc/Dec Drills}

\textbf{Init:} \cee{uint8\_t arr[]={5,7,9,11};} at 0x20000100. \cee{uint8\_t *p = arr;}

\subsection{\cee{x = *p++;}}
\begin{lstlisting}
ldrb r0,[r1],#1
\end{lstlisting}
After: x=5, p=0x101.

\subsection{\cee{x = *++p;}}
\begin{lstlisting}
ldrb r0,[r1,#1]!
\end{lstlisting}
After: x=7, p=0x101.

\subsection{\cee{x = (*p)++;} (return old, increment mem)}
\begin{lstlisting}
ldrb r0,[r1]      @ x = old
add  r2,r0,#1
strb r2,[r1]      @ writeback
\end{lstlisting}
After: x=5, mem[0x100]=6, p unchanged.

\subsection{\cee{x = ++(*p);} (increment mem, return new)}
\begin{lstlisting}
ldrb r0,[r1]
add  r0,r0,#1
strb r0,[r1]
\end{lstlisting}
After: x=6, mem[0x100]=6, p unchanged.

\subsection{\cee{x = ++*p++;}}
``post-inc p, but modify mem at OLD addr, return new value of mem''
\begin{lstlisting}
ldrb r0,[r1],#1   @ x=*old_p, p++
add  r0,r0,#1
strb r0,[r1,#-1]  @ writeback to OLD addr
\end{lstlisting}
After: x=6, mem[0x100]=6, p=0x101.

\subsection{\cee{*q++ = *p++;}}
\begin{lstlisting}
ldrb r0,[r1],#1
strb r0,[r2],#1
\end{lstlisting}

% [TRIM-Part1] Conditional Pre/Post LDR section removed (covered in IT Blocks section)

% ==========================================================
\section{Loops --- Detailed Count}

For an N-element loop with:
\begin{itemize}
\item N iters $\Rightarrow$ N body executions, N+1 cmp+branch checks.
\item Top-tested while: priming-read style $\Rightarrow$ N+1 LDRBs (1 priming + N in loop).
\item while(1)+break $\Rightarrow$ N+1 LDRBs (one per iter incl.\ exit iter).
\end{itemize}

\textbf{Same dynamic LDRB count} between the two; the break-style just removes one source-line of code.

\subsection{do-while (always >= 1 iter)}
\begin{lstlisting}
.lp:
  @ body always runs first
  cmp  r0,r1
  blt  .lp
\end{lstlisting}

\subsection{Counted loop with downto}
\begin{lstlisting}
  mov r4,r5      @ count = N
.lp:
  cbz  r4,.end
  @ body
  sub  r4,r4,#1
  b    .lp
.end:
\end{lstlisting}

\textbf{Trade:} downto uses CBZ (no flags churn) and single SUBS. Up-counting needs CMP each iter.

% [TRIM-Part1] Switch Variants removed (covered in flow-ctrl switch section)

% ==========================================================
\section{C--ASM Pitfall Quiz}

\textbf{Q:} \cee{int8\_t b = 0x80; int32\_t x = b;} --- value of x?\\
\textbf{A:} \textbf{0xFFFFFF80} (sign-extended).

\textbf{Q:} \cee{uint8\_t b = 0x80; int32\_t x = b;} --- value of x?\\
\textbf{A:} \textbf{0x00000080} (zero-extended; the unsigned bool char gets zero-padded).

\textbf{Q:} After \asm{ldr r0,[r1,\#2]!} with $r_1=0x100$, what's $r_1$ after?\\
\textbf{A:} $r_1=0x102$ (pre-index updates BEFORE).

\textbf{Q:} After \asm{ldr r0,[r1],\#2} with $r_1=0x100$, what address was loaded from?\\
\textbf{A:} \textbf{0x100} (post-index updates AFTER).

\textbf{Q:} \asm{push \{r5,r4,r6\}} on SP=0x1020. Where is r6?\\
\textbf{A:} 0x101C (highest addr — regs stored ascending by reg\#).

\textbf{Q:} \cee{p32 + 1} where \cee{p32} is \cee{int32\_t *} = 0x100. Result?\\
\textbf{A:} 0x104 (pointer arith scaled by sizeof(int32\_t)).

\textbf{Q:} \asm{adds r0,r1,r2} where r1=0xFFFFFFFF, r2=0x1. Result and flags?\\
\textbf{A:} r0=0x00000000, Z=1, C=1, N=0, V=0.

\textbf{Q:} CBZ inside IT block --- legal?\\
\textbf{A:} \textbf{No.} CBZ/CBNZ forbidden in IT.

\textbf{Q:} \asm{asr r0,r0,\#1} on r0=0x80000000. Result?\\
\textbf{A:} 0xC0000000 (signed shift, sign bit replicated).

\textbf{Q:} \asm{lsr r0,r0,\#1} on r0=0x80000000. Result?\\
\textbf{A:} 0x40000000 (logical shift, zero fill).

\textbf{Q:} \cee{f((int8\_t)-1)} where \cee{f} takes \cee{int32\_t}. Value in R0?\\
\textbf{A:} 0xFFFFFFFF.

\textbf{Q:} \asm{strb r0,[r1]} where r0=0xDEADBEEF. Memory written?\\
\textbf{A:} 0xEF (only low byte).

% ==========================================================
% [TRIM-Part1] One-liner cheats removed (full versions earlier in Quiz sections)
% [TRIM-Part1] Top-20 patterns moved to Part 2 (duplicated Pattern Catalog)

% ==========================================================
\section{LDR Pseudo Resolution Detail}

\asm{ldr Rd,=expr}: assembler chooses cheapest encoding by value:

\begin{tabular}{@{}lll@{}}
\toprule
Value range & Emitted insn & Notes \\
\midrule
0..0xFF (8-bit) & \asm{movs Rd,\#expr} & 16-bit Thumb \\
0..0xFFFF (16b) & \asm{movw Rd,\#expr} & zeros high 16 \\
big 32-bit      & PC-rel \asm{ldr} + \texttt{.word} & literal pool \\
\bottomrule
\end{tabular}

\textbf{When unsure, use \asm{ldr Rd,=expr}} --- assembler picks the right encoding. Plain \asm{mov Rd,\#imm} only accepts 8-bit (rotated mod-imm) values.

\textbf{Building 32-bit literals manually:}
\begin{lstlisting}
movw r0,#0xBEEF       @ low 16, zero high
movt r0,#0xDEAD       @ high 16, preserve low
@ -> r0 = 0xDEADBEEF
\end{lstlisting}

% ==========================================================
\section{Final Exam Strategy}
\begin{itemize}
\item \textbf{Read the C prototype FIRST.} It dictates LDR variant, signed/unsigned CCS, shift direction.
\item \textbf{Annotate memory layouts.} For LDR traces, write bytes at each address BEFORE running through instructions.
\item \textbf{Track register state explicitly.} After every instruction, write the changed reg's new value.
\item \textbf{Follow the precedence rules.} postfix\,\texttt{++}\, (L1) $\to$ \cee{*} / prefix\,\texttt{++} / cast (L2). Walk through the C carefully.
\item \textbf{Pick addressing mode from inc-order.}
  Prefix-inc $\to$ pre-idx (!); postfix-inc $\to$ post-idx; no inc $\to$ basic offset.
\item \textbf{For IT blocks,} pick the cond as NL; T slots match, E slots inverse. Cond B ONLY last slot. No CBZ in IT.
\item \textbf{For loops,} decide priming vs while(1)-break; CBZ on freshly-loaded bytes.
\item \textbf{Unsigned vs signed.} ALWAYS match: shift (ASR/LSR), CCS (GT/HI), load suf (SH/H).
\item \textbf{Don't forget write-back.} Prefix \cee{++} on a deref value needs STR after the ADD.
\item \textbf{Sanity-check sign-ext.} If MSBit of loaded byte/half = 1, LDRSB/SH produces 0xFFFF-padded high bits.
\item \textbf{Return convention.} $\le$32b $\to$ R0. 64b $\to$ R0:R1. Match return \texttt{bx lr} (leaf) or \texttt{pop \{...,pc\}} (stem).
\end{itemize}

% ==========================================================
\section{Debug Pattern Recognition}

When tracing student-written code, watch for:

\begin{itemize}
\item \textbf{Missing S-suffix} on insn whose flags are tested next $\to$ branch tests stale flags from earlier CMP.
\item \textbf{Wrong CCS family} (e.g.\ BLT with \cee{uint32\_t}) $\to$ wrong direction near 0x80000000.
\item \textbf{Forgotten write-back} after prefix \cee{++} on \cee{*p} $\to$ memory still holds old value.
\item \textbf{LDR offset out of range} for byte/halfword $\to$ assembler error or reg-offset workaround.
\item \textbf{Misaligned LDRD} $\to$ requires 8-byte alignment on some cores; Cortex-M4 forgives but slow.
\item \textbf{IT block too long} $\to$ ITTTTT or 5+ slots $\to$ assembler rejects.
\item \textbf{CBZ in IT} $\to$ assembler rejects.
\item \textbf{Wrong PUSH/POP pairing} $\to$ stack corruption (e.g.\ push \{r4,lr\} but pop \{r4,r5,pc\}).
\item \textbf{Tail call with stem still on stack} $\to$ leaks frame.
\end{itemize}

% ==========================================================
\section{Cortex-M4 Quick Specs}

\begin{tabular}{@{}ll@{}}
\toprule
Feature & Value \\
\midrule
Word size & 32-bit \\
Instr set & Thumb-2 only (no ARM mode) \\
Endianness & little (configurable but typ.\ LE) \\
Pipeline & 3-stage (fetch/decode/execute) \\
Reg file & R0--R15 (16 GP) + APSR/CONTROL \\
SP & MSP (handler) + PSP (thread, RTOS) \\
Bit-banding & SRAM 0x22 region, peripheral 0x42 \\
Vector tbl & 0x08000000 default; relocatable \\
NVIC & 240 IRQs, 0--255 priority \\
SysTick & 24-bit downcounter \\
DSP exts & SIMD, MAC w/ accumulator \\
FPU & VFPv4-SP single-prec (M4F variant) \\
\bottomrule
\end{tabular}

% [MOVE-Part2] Bit-Banding moved to Part 2 (peripheral reference)

% ==========================================================
\section{When to Use What}

\begin{tabular}{@{}ll@{}}
\toprule
Goal & Idiom \\
\midrule
test ==0 in reg & \asm{cbz} (free) or \asm{cmp \#0/beq} \\
test bit set & \asm{tst r,\#mask/bne} \\
clear bit & \asm{bic r,r,\#mask} \\
set bit   & \asm{orr r,r,\#mask} \\
toggle bit & \asm{eor r,r,\#mask} \\
small unsigned mul/div & \asm{lsl}/\asm{lsr} by power-of-2 \\
small signed div & \asm{asr} (signed) \\
3$\cdot$x cheaply & \asm{add r,x,x,lsl \#1} \\
swap upper/lower halves & \asm{ror r,r,\#16} \\
abs value & cmp \#0 / it lt / rsblt \#0 \\
sign mask (-1 or 0) & \asm{asr r,r,\#31} \\
isolate low N bits & \asm{ubfx r,r,\#0,\#N} \\
sign-ext byte & \asm{sxtb r,r} (or LDRSB at load) \\
\bottomrule
\end{tabular}

% ==========================================================
\section{One-Page Summary}

\textbf{Address modes:}
basic \asm{[Ra,\#os]}, pre \asm{[Ra,\#os]!}, post \asm{[Ra],\#os}, reg-off \asm{[Ra,Ri,LSL \#n]}.

\textbf{Sufs:} ---/B/SB/H/SH on LOAD; ---/B/H on STORE.

\textbf{C $\to$ idx:} \cee{*p++}=post, \cee{*++p}=pre, \cee{*(p+k)}=basic, \cee{p[i]}=reg-off.

\textbf{Sign vs unsigned:} CCS GT/HI, shift ASR/LSR, load SH/H. Match!

\textbf{IT:} max 4, T=cond/E=inverse, cond B last only, no CBZ.

\textbf{Loops:} priming-LDRB+CBZ; or while(1)+break.

\textbf{Stem:} push \{lr\}, pop \{pc\}.

\textbf{Modulo:} udiv + mls.

\textbf{64-bit add:} adds + adc; sub: subs + sbc.

\textbf{Sign-ext at LOAD:} LDRSB 0x80$\to$0xFFFFFF80.

\textbf{Sign-ext at CALL:} \cee{int8\_t -1} arg $\to$ R0=0xFFFFFFFF.

\textbf{Q-format:} encode $f \cdot 2^n$ + neg-wrap; decode TC then $/2^n$.

\textbf{PUSH:} reg\# ascending into addr ascending.

\end{multicols*}
\end{document}
